\documentclass[11pt,a4paper,spanish]{book}
\usepackage[a4paper,left=3cm,right=2cm,top=3cm,bottom=2cm]{geometry}
%\usepackage[margin=1in]{geometry}
\usepackage[utf8]{inputenc}
\usepackage[spanish, es-tabla]{babel}
\usepackage[]{enumerate}
\usepackage[]{amssymb}
\usepackage[colorlinks, citecolor=black, filecolor=black, linkcolor=black, urlcolor=black]{hyperref}
\usepackage{graphicx}
\usepackage{float}
\usepackage{multirow}
\usepackage{subfigure}
\usepackage{caption}
\usepackage{minted}
\usepackage[table,xcdraw]{xcolor}
\usepackage{tabularx,ragged2e}
\usepackage{textcomp}
\usepackage{multicol}
\usepackage{tikz}
\usepackage{adjustbox}
%para gramatica bnf
\usepackage{bnf}
%para definir nuevo environment
\usepackage{newfloat}
%para insertar pdf
\usepackage[final]{pdfpages}
%para el arbol de directorios
\usepackage{dirtree}
\usepackage{verbatim} 

\usepackage{algorithm}
\usepackage{algorithmic}

%para el caption de las listas
\DeclareCaptionType{listCap}[Lista]


% declare the floating environment {Grammar}
% this will also define \listofGrammars:
\DeclareFloatingEnvironment[
  % the file extension for the file used to create the list:
  fileext   = logr,% don't use log here!
  % the heading for the list:
  listname  = {List of Grammars},
  % the name used in captions:
  name      = Gramática,
  % the default floating parameters if the environment is used
  % without optional argument:
  placement = htp
]{GrammarEnv}

\newtheorem{definition}{{\bf\sc Definición }}[section]
\newtheorem{theorem}{Teorema}[section]

\newcolumntype{C}{>{\Centering\arraybackslash}X} % centered "X" column

%separacion de columnas multicol
 \setlength{\columnsep}{5mm}

\floatname{algorithm}{Algoritmo}
\renewcommand{\algorithmicrequire}{\textbf{Input:}}
\renewcommand{\algorithmicensure}{\textbf{Output:}}


\begin{document}
	

%Empieza la numeración en números romanos
\frontmatter
% Incluimos la carátula
%Observación: Si bien la página del prefacio dice que sea empty, debería comenzar alli la numeración. Se sugiere numeración romana. Recomenzar la numeración en el primer capítulo de la tesis  con numeración arábiga.

\titlepage

\begin{center}
\ \\
\ \\
\vspace{-1cm}
 

\ \\

\vspace{0.5cm}
{\Large{\bf \sc Universidad Nacional del Comahue}}\\

\ \\
{\Large { \sc Facultad de Informática}}\\

\vspace{-2.5cm}
\mbox{\hspace{-1cm}\includegraphics[width=2.5cm,height=2.5cm]{img/unc.png}\hspace{13cm} \includegraphics[width=2.5cm,height=2.5cm]{img/fai.png}}


\vspace{6cm}

{\Large {\bf\sc Tesis de Licenciado en Ciencias de la Computación}}\\
\ \\
\ \\
{\LARGE {\bf Reestructuración temprana en el proceso de verbalización}}\\
\vspace{3cm}


{\Large Martín Bermudez}\\
\vspace{2cm}

{\Large Esp. Sandra Roger}\\
\ \\
%{\Large [Nombre del CoDirector]}\\

\vfill
{\Large {\sc Neuquén}\hspace{6cm}{\sc Argentina}}\\
\ \\

{\Large 2020}\\

\end{center}

\pagebreak



% Incluimos la página de prefacio.
\ \\
\ \\
\label{pagpref}
\noindent{\LARGE \sc Prefacio}\\
\ \\
\ \\

\ \\

\ \\
\ \\


Esta tesis es presentada como parte de los requisitos finales para optar al grado acad\'emico de 
{\em Licenciado/a en Ciencias de la Computación}, 
otorgado por la Universidad Nacional del Comahue, y no ha sido presentada previamente para la 
obtención de otro título en esta Universidad u otras. La misma es el resultado de la investigación 
llevada a cabo en el Departamento Teoría de la Computación%.... (, el Departamento ... y el 
%Departamento...)
, de la Facultad de Informática, en el período comprendido entre setiembre de
%cuando se aprobó la propuesta
2019 y junio de 2020, bajo la dirección de la Esp. Sandra Roger.%... (y la codirección de ....).




\vspace{3cm}


\ \\
{\flushright Martin Bermudez\\
{\sc Facultad de Informática \\
Universidad Nacional del Comahue}\\
{\em Neuqu\'en, .... de junio de 2020.}\\}

\vfill

\begin{center}
%
\framebox{\begin{minipage}[t]{0.9\columnwidth}%
\begin{flushleft}
\includegraphics[scale=0.035]{img/unc.png}

\vspace{-2cm}
{\large \hspace{5cm}\sc universidad
nacional del comahue} \\
\par\end{flushleft}
\begin{center}
{\large \qquad{}}{ \hspace{2.5cm} Facultad de Informática}
\par\end{center}

\vspace{1cm}

\indent \ \ \ \ \ \ \ \ \ \ \ La presente tesis ha sido aprobada el día ........................., mereciendo la \\
\indent \ \ \ \  \ \ \ \ \ calificación de .............................

\medskip{}

\vspace{1cm}
\end{minipage}}
\end{center}

\pagebreak



% Incluimos la página de dedicatorias. Esta página es opcional.
% Si la página no quiere ser incluida anteponga el símbolo "%" al comienzo de la siguiente línea
\ \\
\ \\
\label{pagdedic}
\noindent{\LARGE \sc Dedicatorias}\\


\vfill
\pagebreak


% Incluimos la página de agradecimientos. Esta página es opcional.
% Si la página no quiere ser incluida anteponga el símbolo "%" al comienzo de la siguiente línea

\ \\
\ \\
\label{pagagrad}
\noindent{\LARGE \sc Agradecimientos}\\
\ \\
\ \\


\vfill
\pagebreak


% Incluimos la página de resumen
\ \\
\ \\
\label{pagresum}
\noindent{\LARGE \sc Resumen}
\\ \\
Las Ontologías juegan un rol clave en la formación de la Web Semántica. Permiten representar el conocimiento de manera formal, brindándole un significado bien definido a los datos y permitiendo que razonadores obtengan información implícita del dominio. 

La Web Semántica es definida como una extensión de la Web actual. Tiene el objetivo de agregar explícitamente una capa de significado a los datos, con el fin de crear un mejor ambiente para poder automatizar tareas complejas. Es en este punto donde las ontologías son útiles, brindando las herramientas adecuadas para estructurar la Web Semántica. 

%Existen varios lenguajes que tienen la capacidad de describir ontologías, como RDF (\emph{Resource Description Framework})~\cite{RDFRecomendW3C} y OWL 2 (\emph{Ontology Web Language})~\cite{RecomendW3C}, siendo este útlimo el recomendado por la La W3C (\emph{World Wide Web Consortium})\footnote{W3C es un consorcio internacional que se encarga de establecer y recomendar estándares para la \emph{World Wide Web}}.
%El  lenguaje  OWL2  (Ontology  Web  Language)~\cite{RecomendW3C}  tiene  base  en  la  lóogica  descriptiva  y  elvocabulario de RDF, lo que lo hace id ́oneo para representar conocimiento y permitir que agentesinform ́aticos puedan razonar sobre ese conocimiento

El lenguaje OWL2 (\emph{Ontology Web Language})~\cite{RecomendW3C}, es una recomendación de la W3C (\emph{World Wide Web Consortium})\footnote{W3C es un consorcio internacional que se encarga de establecer y recomendar estándares para la \emph{World Wide Web}}. Tiene base en la lógica descriptiva y el vocabulario de RDF, lo que lo hace idóneo para representar conocimiento y permitir que agentes informáticos puedan razonar sobre ese conocimiento.

Teniendo en cuenta que las ontologías estructuran los datos de manera formal, son poco comprensibles por usuarios no expertos, y no brindan una cómoda visualización de la información para aquellos que quieran beneficiarse del uso de las tecnologías semánticas.

Si bien existen algunas aplicaciones para  desarrollar y explorar ontologías, éstas no resultan suficientemente satisfactorias para usuarios no expertos o finales. 
Por este motivo, expresar el contenido formal en Lenguaje Natural (LN) resulta atractivo, brindando la capacidad de documentar y expresar ontologías en un lenguaje accesible por usuarios no entrenados en matemáticas o el dominio modelado. 

Para llevar a cabo un sistema que tenga la capacidad de generar lenguaje natural a partir de una ontología, se deben tener en cuenta algunos criterios que impactan sobre el diseño de las ontologías, y sobre los resultados obtenidos de la generación de texto, como dificultad para comprender las relaciones del dominio, oraciones complejas y difíciles de entender, complejidad de la sintaxis y flexibilidad del texto, entre otras. 

Teniendo en cuenta estos criterios, en esta tesis se propone un Sistema de Generación de Texto que mantenga un equilibrio entre la dificultad de integrar el sistema con las ontologías, y la capacidad del sistema de generar texto que resulte útil al usuario final. 

El resultado de este trabajo es un sistema que requiere un mínimo compromiso en el diseño de las ontologías, y produce una salida sintácticamente aceptable y organizada.

Los dos puntos claves en este trabajo son: la \textit{Organización de la Información}, para estructurar el texto con base en las relaciones semánticas, con el fin de reducir la carga cognitiva que exige reconocer las relaciones del dominio; y la \textit{Generación del Texto en Lenguaje Natural}, maximizando la cohesión de las oraciones.

\vfill
\pagebreak


% Incluimos la página de resumen en inglés
\ \\
\ \\
\label{pagsumm}
\noindent{\LARGE \sc Abstract}
\\ \\
Ontologies play a key role in the formation of the Semantic Web. They allow the representation of knowledge in a formal way, giving a well-defined meaning to the data and allowing reasoners to obtain implicit information from the domain.

The Semantic Web is defined as an extension of the current Web. It aims to explicitly add a layer of meaning to the data, in order to create a better environment for automating complex tasks. It is at this point that ontologies are useful, providing the appropriate tools to structure the Semantic Web.

Taking into account that ontologies structure data in a formal way, they are poorly understood by non-expert users, and do not provide a comfortable display of information for those who want to benefit from the use of semantic technologies.

Although there are some applications for developing and exploring ontologies, they are not satisfactory enough for non-expert or end users.
For this reason, expressing formal content in Natural Language (LN) is attractive, providing the ability to document and express ontologies in a language accessible by users not trained in mathematics or the modeled domain.


However, expressing the formal content in LN is not useful enough if only isolated sentences are generated. For a text to be beneficial to an end user, it must be organized and consistent. In this way the relationships between the concepts can be captured, to understand the complete domain and not only interpret its isolated axioms. For this reason, a Natural Language Generation System (SGLN) was designed and developed, which allows an organized text to be obtained from the content of an ontology in OWL 2 language.


The result of this work is a system that requires a minimum commitment in the design of ontologies, and produces a syntactically acceptable and organized output.

The two key points in this work are: the \textit{Organization of Information}, to structure the text based on semantic relationships, in order to reduce the cognitive load that requires recognizing domain relationships; and the \textit{Natural Language Text Generation}, maximizing the cohesion of sentences.


\vfill
\pagebreak



% Tabla de contenido o Indice del contenido
\tableofcontents{}

% Lista de figuras o Indice de figuras
\listoffigures

% Lista de tablas o Indice de Tablas
%\listoftables

\newpage

%Empieza la numeración en números arábigos
\mainmatter



\chapter{Introducción}
La Web Semántica puede ser vista como una colección de estándares y tecnologías que permiten a las máquinas entender el significado (semántica) de la información disponible en la Web~\cite{yu2011developer}. Tiene como objetivo enriquecer la web actual agregando explícitamente una capa de significado a los datos. De esta manera, se crea un mejor ambiente para poder automatizar tareas complejas que actualmente se realizan por los humanos de forma manual. Además, crea un espacio donde la información tendría un significado bien definido, de manera que pudiera ser interpretada tanto por agentes humanos como por agentes computarizados.

%MAYO el owl  no es solo para describir sino  para razonar tambien. Para representar tenemos tambien a RDF por ejpemplo. Habría que arreglar. AGREGADO RDF Y RAZONAMIENTO
Para describir el dominio de interés en la Web Semántica, se adoptó el uso de ontologías. Algunos lenguajes conocidos para representar ontologías son RDF (\emph{Resource Description Framework})~\cite{RDFRecomendW3C} y OWL 2 (\emph{Ontology Web Language}). La W3C (\emph{World Wide Web Consortium })\footnote{W3C es un consorcio internacional que se encarga de establecer y recomendar estándares para la \emph{World Wide Web}} recomienda la representación de la información en la Web Semántica con el lenguaje OWL 2~\cite{RecomendW3C}. El lenguaje OWL 2 tiene base en la lógica descriptiva y el vocabulario de RDF, lo que lo hace idóneo para representar conocimiento y permitir que agentes informáticos puedan razonar sobre ese conocimiento.

Como las ontologías estructuran los datos de manera formal, son poco comprensibles por usuarios no expertos, y no brindan una cómoda visualización de la información para aquellos que quieran beneficiarse del uso de las tecnologías semánticas.

Actualmente existen algunas aplicaciones para  desarrollar y explorar ontologías, como Pro-tégé~\cite{protege}, Mindswap~\cite{golbeck2002new}, OntoEdit~\cite{sure2002ontoedit}, OntoLingua~\cite{farquhar1997ontolingua}, WebOnto~\cite{domingue1998tadzebao}, entre otras, pero no resultan suficientemente satisfactorias para usuarios no expertos, o usuarios finales. 

Por este motivo, expresar el contenido formal en Lenguaje Natural (LN) resulta atractivo, brindando la capacidad de documentar y expresar ontologías en un lenguaje accesible por usuarios no entrenados en matemáticas o el dominio modelado. Un Sistema de Generación de Lenguaje Natural (SGLN) puede tener diferentes objetivos en cuanto a la salida que producen, ya sea enfocarse en la facilidad de lectura, la capacidad de poder generar diferentes oraciones para expresar la misma información,  enfocarse en la organización del contenido, entre otras. En el ámbito de las ontologías existen algunos trabajos que involucran el procesamiento de lenguaje natural~\cite{moreno2018ontologia}\cite{perez2002explotacion}\cite{vallez2009web}. Como se verá en la Sección~\ref{sec:trabajos_rel}, cada enfoque adoptado por los trabajos tiene ventajas y desventajas, por ejemplo en cuanto a la integración con las ontologías, debido a la necesidad de modificar la ontología subyacente o de requerir convenciones de nombrado; como en su potencial para resultarle útil al usuario final, debido a la complejidad de las oraciones, la falta de flexibilidad léxica, y la organización del contenido del texto, lo que resulta en una pobre comprensión el texto.  

Teniendo en cuenta estas ventajas y desventajas, en esta tesis se propone un Sistema de Generación de Texto que mantenga un equilibrio entre la dificultad de integrar el sistema con las ontologías, y la capacidad del sistema de generar texto que resulte útil al usuario final. El Sistema de Generación de Texto debe   ser capaz de recibir una ontología de la Web Semántica y generar un documento de texto organizado y en lenguaje natural. No será necesario agregar información a la base de conocimiento, mientras que se espera alcanzar una salida que sea sintácticamente aceptable.

Los dos puntos claves en este trabajo son: la \textit{Organización de la Información}, para estructurar el texto con base en las relaciones semánticas, con el fin de reducir la carga cognitiva que exige reconocer las relaciones del dominio; y la \textit{Generación del Texto en Lenguaje Natural}, maximizando la cohesión de las oraciones.

\section{Objetivos}\label{Intro:objetivo}
El objetivo de este trabajo es diseñar e implementar una aplicación capaz de procesar el contenido de una ontología y generar un texto en lenguaje natural. El texto resultante debe estar organizado y  expresado con una sintaxis aceptable, para que sea accesible a cualquier usuario cuyo interés sea el dominio modelado, pero no tenga las herramientas necesarias para comprender las lógicas subyacentes.

Para alcanzar tales objetivos, se pretende desglosarlo en dos sub-objetivos, que se corresponden con sus dos problemas principales: organizar y verbalizar la información. Para el primer problema se propone un preprocesamiento de la base de conocimiento, con el fin de analizar su estructura y establecer una jerarquía de tópicos. Luego, se abordará el problema de la verbalización, el cual tendrá como resultado el texto final, cuya organización será establecida por la jerarquía de tópicos.

Respecto al equilibrio entre la modificación de la base de conocimiento y la calidad del texto de salida, queremos evitar modificar y agregar información ajena al dominio modelado, pidiendo un compromiso mínimo en la información que provee la ontología, para realizar un Sistema de Generación de Texto lo más independiente y genérico posible.

\section{Motivación}
La motivación principal es facilitar la comprensión de ontologías. Utilizar procesos de generación de texto automáticos resultan convenientes para generar documentos de texto que sean comprensibles por cualquier usuario.

%Por otro lado, también queremos hacer uso de teorías lingüísticas, que se enfocan tanto en la organización del texto como en la organización del contenido de sus sentencias, es decir, se enfocan en los diferentes niveles de coherencia del texto. Esperamos poder utilizar estos conceptos para motivar el desarrollo de sistemas basados en teorías lingüísticas, y no tanto en el desarrollo de inteligencias artificiales basadas en algoritmos de aprendizaje, con el fin de acercar más las disciplinas tales como Lingüística Teórica y Ciencias de la Computación. 
%MAYO yo sacaría este párrafo siguiente .. queda como que falta algo más ...  lo dí vuelta fijate .. 
%-LO COMENTE ASI NO APARECE EL PARRAFO-.
%El papel que ha desempeñado la Lingüística Teórica en el Procesamiento de Lenguaje Natural ha sido poco notorio~\cite{perinan2012defensa}.
%Por otro lado, también queremos hacer uso de teorías lingüísticas, que se enfocan tanto en la organización del texto como en la organización del contenido de sus sentencias, es decir, se enfocan en los diferentes niveles de coherencia del texto. Dado que el papel que ha desempeñado la Lingüística Teórica en el Procesamiento de Lenguaje Natural ha sido poco notorio~\cite{perinan2012defensa}, esperamos poder utilizar estos conceptos para motivar el desarrollo de sistemas basados en estas teorías lingüísticas, y no tanto en el desarrollo de inteligencias artificiales basadas en algoritmos de aprendizaje, con el fin de acercar más las disciplinas tales como Lingüística Teórica y Ciencias de la Computación. 


\chapter{Antecedentes conceptuales}

\section{Trabajos relacionados}
%%copiado y pegado de la propuesta
Varios trabajos involucran el procesamiento de lenguaje natural en el ámbito de las ontologías~\cite{moreno2018ontologia}\cite{perez2002explotacion}\cite{vallez2009web}. Por un lado, existen herramientas que se enfocan en la verbalización de un conjunto de axiomas pertenecientes a una ontología, con el fin de generar un conjunto de sentencias que puedan ser utilizadas para reconocer la información perteneciente al dominio modelado, o documentar la ontología, utilizando lenguaje natural. En este tipo de trabajos se encuentran algunas herramientas como OntoVerbal~\cite{liang2013ontoverbal}, que utiliza la frecuencia con la que se utilizan conjuntos de axiomas para la descripción de entidades, y describe patrones lingüísticos para verbalizar los conjuntos de mayor frecuencia. Otra característica de OntoVerbal es que presenta en un mismo párrafo las sentencias que están estrechamente relacionadas. %Además, agrupa los axiomas alrededor de un tópico para crear párrafos. 
%Además, genera la creación de párrafos mediante la agrupación de axiomas alrededor de un tópico dado.

Por su parte, NaturalOWL~\cite{galanis2007generating}, se enfoca en la flexibilidad y fluidez del texto, agregando información lingüística a la ontología para verbalizar los axiomas.

Por otro lado, se han desarrollado herramientas para asistir durante el desarrollo de la ontología, permitiendo acceder a la información modelada en lógicas descriptivas en un formato de texto, y poder editarlo en lenguaje natural. Entre las herramientas desarrolladas se encuentran CLOnE~\cite{power2010complexity} y Aqualog~\cite{lopez2005aqualog},completamente orientadas a la comprensión de texto y no a la generación; y Attempto Tools~\cite{attempto}, que utiliza Attempto Controlled English  (ACE) que es un lenguaje natural controlado ~\cite{CNL}, para servir como lenguaje de representación de conocimiento, y de esta manera poder utilizar representaciones formales, sin necesidad de comprender la lógica interna. La ventaja de usar ACE radica en la posibilidad de generar lenguaje natural desde una representación de conocimiento en lógica, y viceversa.

Si bien estas herramientas cumplen con la verbalización de axiomas y asisten en el uso de ontologías evitando que los usuarios deban comprender las lógicas descriptivas y la sintaxis OWL, la comprensión de cada axioma aislado no asegura la comprensión del dominio completo. Un factor importante en la comprensión del dominio es la capacidad de relacionar los temas de manera coherente. La verbalización de axiomas es necesaria pero no suficiente para comprender las temáticas modeladas en la ontología. Existen conjuntos de axiomas que están relacionados y son complementarios para la descripción de entidades, y que expresados en conjunto como una unidad de información textual transmiten mejor la intención de ese conjunto de axiomas. OntoVerbal es la herramienta que tiene en cuenta este criterio, sin embargo, no planea una estructura de documento, por lo que no garantiza que el orden en el que aparezcan los párrafos maximice la coherencia.

En NaturalOWL se tiene en cuenta la calidad del texto de salida, pero compromete la representación del dominio, agregando información lingüística y ensuciando el dominio modelado. Por último en Attempto Tools se busca un equilibrio en la calidad del texto sin comprometer el dominio modelado, pero no se hace énfasis en la comprensión del dominio, ya que únicamente se traducen los axiomas en sentencias aisladas. 

Por estos motivo, en este trabajo se propone diseñar un algoritmo que agrupe la información relacionada a través de los temas principales presentes en la ontología. En este sentido, se propone una técnica alternativa a la de Reiter y Dale~\cite{reiter1997building} quienes realizan todo el ordenamiento de las sentencias e información involucradas en las últimas tareas de las verbalización. Por el contrario, en nuestra propuesta se pretende organizar la mayor parte de la información antes del proceso de verbalización en sí mismo, por considerarse que en éste punto es donde se puede lograr una manipulación más fácil y limpia de la información y dejando solamente el ordenamiento mínimo en etapas posteriores. De esta manera, esperamos que el resultado de la verbalización sean sentencias en lenguaje natural relacionadas semánticamente, que comprendan unidades textuales de mayor granularidad y no únicamente unidades textuales  aisladas (y posiblemente desordenadas) a nivel de oración  que cumplan una determinada sintaxis.

Con esto en mente, se propone diseñar e implementar una herramienta genérica para elegir, ordenar y generar un texto en lenguaje natural a partir de una ontología en OWL, con el fin de facilitar el entendimiento del dominio modelado, y asistir durante el proceso de modelado. Este trabajo hará foco principalmente en generación de texto en lenguaje español, y presentará un prototipo en lenguaje inglés.

\section{Ontologías}
Las ontologías definen conceptos y relaciones de algún dominio, de forma compartida y consensuada; y  esta conceptualización debe ser representada de una manera formal, legible y utilizable por los ordenadores~\cite{tello2001ontologias}. 

Las ontologías se componen de~\cite{llop2007especificacion}:
\begin{itemize}
    \item Clases: representan un conjunto de objetos.
    \item Instancias: cada instancia representa un objeto específico, que pertenece a una única clase.
    \item Propiedades: definen las propiedades que pueden tener los objetos.
    \item Relaciones: representan las relaciones que pueden haber entre los objetos.
\end{itemize}
Las ontologías para la Web Semántica, adicionalmente consisten  de reglas de inferencia~\cite{SWBernersLee}, que permiten obtener conocimiento implícito para la toma de decisiones.

Dado que deben estar representadas formalmente, se expresan en lenguajes formales, tal como OWL.

\subsection{Lenguaje OWL}
El lenguaje OWL (\emph{Web Ontology Language}) es un lenguaje ontológico basado en lógicas descriptivas y RDF. 

Tiene una sintaxis robusta, con constructores y axiomas que permiten definir Clases.

Como existen varios tipos de Lógicas Descriptivas, la sintaxis OWL permite definir tres niveles de expresividad, limitando los constructores usados en cada nivel, para tener control sobre la complejidad computacional a la hora de inferir conocimiento~\cite{tello2001ontologias}. Las tres sintaxis son OWL Lite, OWL DL y OWL Full, de menor a mayor expresividad respectivamente. Los constructores y axiomas disponibles son: 
\begin{multicols}{3}
\begin{itemize}
    \item owl:Class
    \item rdfs:subClassOf
    \item rdfs:subPropertyOf
    \item rdfs:domain
    \item rdfs:range
    \item owl:Individual 
    \item owl:ObjectProperty
    \item owl:DatatypeProperty
    \item rdfs:subClassOf
    \item owl:intersectionOf
    \item owl:minCardinality
    \item owl:maxCardinality
    \item owl:Cardinality
    \item owl:complementOf
    \item owl:unionOf
    \item owl:equivalentClass
    \item owl:disjointWith
    \item owl:oneOf
    \item owl:hasValue
    \item owl:AllDifferent
    \item owl:allValuesFrom
    \item owl:someValuesFrom 
    \item owl:disjointWith
    \item owl:inverseOf
\end{itemize}
\end{multicols}

 
\section{Generación de lenguaje natural}
\label{sec:tareas_gnl}
La generación de lenguaje natural es un proceso que tiene como principal objetivo la generación de texto a partir de datos de entrada.
Durante el proceso de generación de lenguaje natural, la entrada atraviesa diferentes etapas, a través de las cuales se modifica, hasta alcanzar la última etapa en la que se obtiene un el resultado en lenguaje natural.

Según el tipo de entrada, un Sistema de Generación de Lenguaje Natural (SGLN) puede clasificarse en Datos a Texto o Texto a Texto~\cite{vicente2015generacion}.
\begin{itemize}
    \item Datos a Texto: la entrada es un conjunto de datos estructurados. Por ejemplo, pueden ser datos numéricos, tal como la descripción del clima; datos  tabulares~\cite{mahapatra2016statistical}; datos en formato JSON\footnote{referencia a tesis matias?}; o estructuras más complejas como ontologías.
    \item Texto a Texto: la entrada puede ser texto u oraciones aisladas. Para trabajar con este tipo de entrada, se utilizan enfoques estadísticos~\cite{mittal1999ultra}; la transformación del texto a estructuras de datos~\cite{saldanha2004creation}; o modelos formales~\cite{guo2018long}.
\end{itemize}

\subsection{Tareas de la generación de lenguaje natural}
En el ámbito de la Generación de Lenguaje Natural, se engloban las tareas de cada nivel de abstracción en una etapa particular. Existen tres etapas principales: \emph{macro planning}, \emph{micro planning} y \emph{realización gramatical}~\cite{vicente2015generacion}.

\begin{itemize}
    \item  \emph{Macro planning:} agrupa las tareas de nivel más abstracto. Se encarga de seleccionar la información que se usará en la generación de texto y de organizarla a nivel de capítulos, secciones, párrafos, y sentencias.
    \item \emph{Micro planning:} es la etapa asociada a la construcción de sentencias, agrupa las tareas de nivel de abstracción intermedio. Dado un conjunto de datos que deben estar presentes en una sentencia, el \emph{micro planning} decide en qué orden aparecerán y cómo se combinarán para producir una oración cohesiva. 
    Para esto, se selecciona las configuraciones léxicas de las palabras, se elige qué entidades se reemplazarán por una expresión de referencia, y qué oraciones se pueden combinar para mejorar la fluidez del texto.
    \item \emph{Realización gramatical y expresiones de referencia:} es la última etapa, donde se seleccionan las palabras adecuadas según lo establecido en el \emph{micro planning}. Puede involucrar, por ejemplo, el uso de algoritmos de conjugación de verbos. Expresiones de referencia se refiere a la referenciación de entidades a través de expresiones que las identifiquen.
\end{itemize}

\subsection{Arquitecturas de SGLN}
%https://pdfs.semanticscholar.org/bfc1/44401be9cb93493b0974afb561b8d3c2d1ec.pdf
La forma en que se comunican las tareas descritas en la sección anterior definen la arquitectura del sistema~\cite{ibanez2004arquitectura}. La arquitectura más generalizada es el \emph{pipeline}~\cite{hervas2008descripcion}, representado en la Figura~\ref{fig:arq_pipeline}.
%https://pdfs.semanticscholar.org/6cd2/54a9a4ea49aee61c3add58c11c29162428ae.pdf
Esta arquitectura tiene la ventaja de ser simple de implementar, pero es poco flexible, por lo que las decisiones tomadas en etapas tempranas deben mantenerse a lo largo del pipeline sin posibilidad de mejora.

\begin{figure}
    \centering
    \includegraphics[]{img/antecedentes/pipeline.pdf}
    \caption{Arquitectura pipeline}
    \label{fig:arq_pipeline}
\end{figure}
 

Cada arquitectura se enfoca en dos aspectos fundamentales: la independencia de los módulos (qué tan acoplado/desacoplado está el sistema) y el flujo de control entre ellos. Desde el punto de vista de la separación de los módulos, se pueden encontrar arquitecturas integradas (o monolíticas) donde todo el sistema está formado por un único módulo, hasta sistemas en los que a cada tarea se le asocia un módulo separado~\cite{hervas2008descripcion}. Desde el punto de vista del flujo de control, se encuentran arquitecturas como pipeline donde la información fluye en un único sentido a través de los módulos, hasta arquitecturas de pizarra, donde la información se escribe en un sitio en común a todos los módulos.
Una descripción más detallada de cada arquitectura se encuentra  en~\cite{ibanez2004arquitectura}.

\subsection{Enfoque para abordar la GNL: coherencia local y global}
En la generación de texto automática, además de desear cumplir ciertas características léxicas y sintácticas del lenguaje destino, se requiere que el texto sea coherente. Para alcanzar esta coherencia textual, se puede dividir el estudio de coherencia en dos~\cite{van2005estructuras}: la coherencia local, asociada a la conexión semántica entre oraciones; y la coherencia global, asociada al asunto del texto.

Para abordar el problema de la coherencia local, nos basaremos en lo postulado por la Teoría de Centrado~\cite{grosz1986towards}.
En cuanto a la coherencia global, utilizaremos como recurso a la Teoría de Veins~\cite{cristea1998veins}.

\subsubsection{Teoría de Centrado}
La Teoría de Centrado es una teoría de coherencia local y tópicos de discurso. Por un lado, intenta caracterizar los discursos coherentes entre entidades, teniendo en cuenta la forma en la que se introducen y discuten sus entidades. Por otro lado, se encarga de predecir qué entidades serán más destacadas en un momento dado~\cite{poesio2004centering}.

Esta teoría propone un conjunto de reglas, que de aplicarlas favorece la coherencia del discurso. La afirmación más fundamental de la Teoría de Centrado es que, en la medida en que un discurso se adhiera a las restricciones centrales, su coherencia aumentará y la carga de inferencia sobre el oyente disminuirá~\cite{kruijff1997topics}. Esta característica que aclama, resulta conveniente para el objetivo que busca nuestro trabajo, respecto a minimizar los esfuerzos de un usuario para comprender el dominio modelado en una ontología.

Según~\cite{vicente2015generacion}, ``Un elemento de una parte del discurso a nivel local se constituye como el foco de atención o centro de ese contexto, como la entidad más relevante a la que se refiere el resto de proferencias\footnote{Se entiende como proferencia al conjunto de técnicas o herramientas que nos permiten adelantar acontecimientos o hechos futuros a partir de datos del pasado}''. Desde el punto de vista de la lingüística computacional, esto resulta útil para abordar el problema de la anáfora\footnote{Una anáfora es una figura retórica que consiste en la repetición de una o varias palabras al principio de un verso o enunciado.}. 

\subsubsection{Teoría de Veins}
La Teoría de Veins es una extensión de la Teoría de Centrado. Esta teoría propone un modelo para la descripción de la coherencia global. La estructura jerárquica del texto permite la relación (a través de referencias) entre unidades textuales aún cuando no se encuentran en una disposición lineal~\cite{cristea2005motivations}, garantizando que se mantiene la coherencia del texto.

\section{Teoría de grafos}
La teoría de grafos se encarga de estudiar las propiedades de los grafos. Un grafo es una estructura de datos que consiste de un conjunto de vértices conectados por un conjunto de arcos que pueden ser usados para modelar relaciones entre los objetos en una colección~\cite{mihalcea2011graph}.

Algunos aspectos que se estudian de los grafos son su estructura (jerárquica, árbol, conexos, bipartitos, etc.), las medidas de sus nodos (centralidad relativa, absoluta~\cite{freeman1978centrality}), los algoritmos de recorridos (profundidad, anchura, etc.), entre otros.

La teoría de grafos tiene diversas aplicaciones, por ejemplo el análisis de redes sociales, representación de mapas conceptuales, diseño de circuitos eléctricos, topologías de redes de computadoras, etc. Actualmente se estudia la teoría de grafos y el procesamiento de lenguaje natural en conjunto, para hallar solución a diferentes problemas del procesamiento de texto, como el etiquetado léxico, clasificación de texto y clustering~\cite{mihalcea2011graph}, desambiguación semántica de palabras~\cite{MihalceaSinhaDisambiguation}\cite{ArabJahromiDisambiguation}, extracción de keywords~\cite{Litvak:2008:GKE:1613172.1613178}\cite{beliga2015overview}, generación de resúmenes~\cite{plaza2011uso}, entre otras.

\subsection{Criterio de centralidad de un nodo}
El criterio de centralidad es alguna función del grado de un nodo~\cite{freeman1978centrality}. El grado de un nodo $n$ es la cuenta de la cantidad de nodos adyacentes a $n$. 

La medida de centralidad es utilizada para cuantificar la importancia de un nodo~\cite{2018transfer}. Un ejemplo del cálculo de centralidad es \emph{PageRank}~\cite{page1999pagerank} el algoritmo utilizado por Google para calcular la relevancia de los documentos web.

Existen diferentes medidas de centralidad, algunas de ellas son:
\begin{itemize}
    \item \emph{Degree Centrality}: en grafos no dirigidos es la cuenta de la cantidad de conexiones de un nodo; en un grafo dirigido, se tiene en cuenta la dirección de la conexión, siendo outdegree la medida de la cantidad de conexiones de salida de un nodo, e indegree la medida de la cantidad de conexiones entrantes al nodo. 
    \item \emph{Closeness Centrality}: es la longitud promedio del camino más corto entre el nodo y todos los demás nodos en el grafo.
    \item \emph{Between Centrality}: calcula la cantidad de veces que un nodo actúa como un puente en el camino más corto entre otros dos nodos.
\end{itemize}

\subsection{Algoritmos de recorridos de grafos}
Se utilizan para recorrer los nodos de un grafo. Dos clasificaciones principales:
\begin{itemize}
    \item Recorrido en profundidad: dado un nodo $n$, se visita un nodo $n'$ adyacente, explorando cada camino que salga de él. Hasta que no se haya finalizado de explorar uno de los caminos no se comienza con el siguiente. Un camino deja de explorarse cuando se llega a un nodo ya visitado. 
    \item Recorrido en anchura: dado un nodo $n$, se visitan todos sus nodos adyacentes. Sea $n'$ un adyacente de $n$, no se visita ningún adyacente de $n'$ hasta que ya se hayan visitado todos los adyacentes de $n$. 
\end{itemize}

\section{Procesamiento de lenguaje natural y teoría de grafos}
La representación de texto como una estructura de datos basada en grafos ha sido usada para dar soporte a las tareas de procesamiento de lenguaje natural. En TextGraph Workshop Series\footnote{http://www.textgraphs.org} se pone a disposición una plataforma para el intercambio de ideas que beneficien ambos campos de investigación. 

La representación de datos como un grafo beneficia el PNL poniendo a disposición los algoritmos conocidos, permitiendo la navegación del grafo para extraer y clasificar el contenido. Algunas tareas de PLN que son resueltas utilizando grafos son la creación de resúmenes~\cite{mihalcea2004graph}\cite{erkan2004lexrank} y minería de texto~\cite{Jin:2007:GTR:1244002.1244182}\cite{beliga2015overview}.

\begin{itemize}
    \item Resúmenes: ``La generación automática de resúmenes consiste en la creación de una versión reducida de uno o varios documentos por parte de un programa de ordenador, de tal forma que el resumen producido condense la información importante del texto de entrada''~\cite{plaza2011uso}.
    \item Minería de texto: ``Es el proceso de extracción automática de información fundamental de textos,  detección  automática  de temas  predominantes en un conjunto de documentos y búsqueda de textos relevantes mediante consultas de grandes prestaciones y flexibilidad''~\cite{brun2004articulo}.
\end{itemize}

\section{Ontologías y teoría de grafos}
Tanto en la academia como en la industria~\cite{soylu2018navigating} se han utilizado grafos para la representación de conocimiento. Este enfoque permite explotar las propiedades de un grafo dentro del área de la visualización de la información, por ejemplo para la representación y navegación de ontologías~\cite{escarza2005visualizacion}.

El uso de herencia múltiple en ontologías genera estructuras de datos en forma de grafos dirigidos, sin embargo se pueden usar diferentes criterios para mostrar el contenido de una ontología, por ejemplo Protégé utiliza una visualización en forma de árbol para mostrar la jerarquía de clases, donde cada nodo del árbol es una clase~\cite{protegehierarchy}. Como es posible que un nodo sea hijo de más de un padre, existe redundancia en la visualización, ya que un nodo puede repetirse en dos ramas diferentes.

En~\cite{mellish2008natural} presentan una investigación acerca de \emph{Natural Language Directed Inference from Ontologies} (NLDI), en el cual buscan determinar el contenido para describir una entidad presente en una ontología. Los axiomas elegidos para tal descripción surgen de aplicar un algoritmo \emph{best-fist search} sobre un grafo generado a partir de los axiomas de la ontología.

En~\cite{crampes2007ontology} comparan dos aplicaciones que utilizan ontologías para soportar y dirigir la navegación conceptual. En ellas, aplican criterios de selección y ordenamientos de conceptos a partir de las relaciones de las entidades en la ontología del dominio y de reglas específicas de una ontología pedagógica. 

De estos trabajos se observan dos aspectos importantes para integrar la teoría de grafos y la representación de las ontologías:
\begin{itemize}
    \item Determinar qué información de la ontología será usada para crear los nodos y los enlaces del grafo.
    \item Establecer un método de navegación del grafo.
\end{itemize}

Dado que se trabaja con grafos semánticos, en el contexto de la generación de lenguaje natural es deseable aprovechar al máximo la semántica de las relaciones, ya que cada ruta entre los axiomas corresponden a diferentes transiciones posibles en un texto coherente~\cite{mellish2008natural}.


\section{Web Semántica y Teoría de Redes}
%https://riunet.upv.es/bitstream/handle/10251/86297/SOLARES%20-%20Redes%20aleatorias%2C%20de%20peque%C3%B1o%20mundo%20y%20libres%20de%20escala.pdf?sequence=1
La Teoría de las Redes surge de la Teoría de Grafos, enfocándose en el estudio de la estructura topológica del grafo~\cite{solares2017redes}. Dentro del estudio de las redes, se distinguen las redes aleatorias, redes de pequeños mundos y las redes libres de escala. Estas últimas son de nuestro interés, ya que las ontologías de la Web Semántica se caracterizan por ser redes libres de escala~\cite{zhang2008scale}.

En una red libre de escala, ``la distribución de los grados de la red obedece a una ley de potencia, debido a la presencia de ciertos tipos de nodos llamados \emph{hubs} que se caracterizan por estar conectados a un gran número de otros nodos, es decir, que el grado de estos nodos es muy grande en relación a los demás nodos de toda la red''~\cite{solares2017redes}.

%\chapter{Presentación del problema}

Para cumplir con el objetivo propuesto, se dividirá el desarrollo del sistema de generación de texto en dos problemas. El primero consiste en crear un módulo que preprocese la entrada para organizar y estructurar la información de forma más significativa. Reiter E. ha propuesto esta etapa previa de preprocesamiento~\cite{reiter2007architecture} para el análisis y la interpretación de datos, aunque fue aplicada en una entrada de datos en bruto en lugar de una base de conocimiento.

El segundo problema consiste en diseñar e implementar el módulo que toma como entrada la información organizada y la convierte en un documento de texto. Para esto usaremos como referencia las tareas descritas en la Sección~\ref{sec:tareas_gnl}.

Se puede apreciar la composición de nuestro sistema de generación de lenguaje natural en la Figura~\ref{fig:modulos_sgln}.

\begin{figure}
    \centering
    \includegraphics[width=12cm, height=4cm]{img/presentacion_problema/modulos_sgln.pdf}
    \caption{Módulos que componen el sistema de generación de lenguaje natural}
    \label{fig:modulos_sgln}
\end{figure}

\section{Presentación del problema de la organización de la información para la coherencia del texto}
\label{sec:problema_coherencia-texto}
 En la generación de texto automática, además de desear cumplir ciertas características léxicas y sintácticas del lenguaje destino, se requiere que las oraciones del texto posean una conexión semántica, para dar pie a la coherencia del texto. 
 %%%
Si bien existen diferentes formas de organizar la información para formar un texto coherente, debemos exigir que entre las unidades de comunicación adyacentes, haya relación entre los tópicos que presentan, y no únicamente una relación semántica. Por ejemplo, el siguiente fragmento de texto fue presentado en \cite{van1983ciencia}:
``Compré esta máquina de escribir en Nueva York. Nueva York es una gran ciudad de USA. Las grandes ciudades a veces tienen serios problemas financieros''. A medida que se avanza en el texto, se presenta una conexión lineal entre algunos componentes de cada oración, pero no parece transmitir un mensaje claro, pues, los tópicos de cada oración son distintos. De igual manera sucede entre párrafos y capítulos que no mantienen un tópico en común. 

Este salto entre diferentes temas dificulta la interpretación del texto. Considerando el ámbito de las ontologías, resulta conveniente mantener la centralidad del tema mientras se describe el dominio, para que el lector pueda hacer las conexiones adecuadas entre las entidades del dominio.

Para encarar esta problemática, existen diferentes investigaciones que buscan organizar sentencias midiendo la fuerza de sus relaciones [referenciar a alguna de las investigaciones en macroplanning]. 
\\

En una ontología, la información se representa en forma de entidades relacionadas con otras entidades. Una manera de representar visualmente un conjunto de entidades, sus relaciones y atributos es mediante el uso de grafos \ref{escarza2005visualizacion}. Podemos valernos de la estructura del grafo, para medir el grado de relación entre las entidades, antes del proceso de verbalización. Se pueden utilizar diferentes criterios para medir el grado de relación. A continuación veremos un ejemplo a partir de la Figura~\ref{fig:pizza.owl}, la cual contiene representada gráficamente un subconjunto de la ontología pizza.owl\footnote{https://protege.stanford.edu/ontologies/pizza/pizza.owl}.

\begin{figure}
\centering
\subfigure[Jerarquía de clases de la ontología pizza]{\includegraphics[width=70mm]{img/presentacion_problema/onto_pizza.pdf}}
\subfigure[Jerarquía de propiedades de la ontología pizza]{\includegraphics[width=3cm]{img/presentacion_problema/onto_pizza_properties.pdf}}
\caption{Subconjunto de la ontología de pizzas.} \label{fig:pizza.owl}
\end{figure}

Aún antes de generar las proposiciones que reflejan la información del grafo, se puede ver que hay formas más adecuadas que otras para comenzar a transmitirle a un receptor toda la información del grafo. Algunos de los posibles tópicos que engloban o relacionan a la mayoría de los datos expuestos son ``Tipo de comidas'', o ``Ingredientes de la pizza'', siendo que, la mayor cantidad de información está relacionada a la pizza. 

Para identificar los nodos más relevantes, pueden usarse diferentes criterios sobre la estructura del grafo, tal como la disposición jerárquica, o la cantidad de relaciones (el grado) de cada nodo. También se puede hacer uso de la jerarquía de las propiedades, para reconocer su tipo más general, como en \emph{tiene cobertura} y \emph{tiene base}, con el fin de referirse a ellas como ``Ingredientes''.

Para mostrar un ejemplo, supondremos una ontología que represente el texto de la compra de la máquina de escribir. Puede verse el grafo de la ontología en la Figura~\ref{fig:maquina_escr}.

\begin{figure}
    \centering
    \includegraphics[scale=0.65]{img/presentacion_problema/onto_maq_escr.pdf}
    \caption{Ontología del texto de la máquina de escribir.}
    \label{fig:maquina_escr}
\end{figure}

Basándonos en el texto propuesto, puede verse que para reproducirlo a partir de la ontología propuesta, se debe recorrer el grafo en profundidad, comenzando desde la clase \emph{Persona}. De esta manera, aunque los datos tengan una relación semántica, en la generación del texto no queda claro el objetivo o la finalidad del texto. 

Si en lugar de recorrer el grafo en profundidad, se recorre en anchura, se produce un texto con más sentido, ya que centra como tópico principal a la máquina de escribir. Por ejemplo, si comenzamos desde \emph{Persona}, podemos formar el texto ``Compré una máquina de escribir. Es de marca X y modelo Y, con color negro. Fue fabricada en New York''.


En estos ejemplos hemos presentado dos problemáticas a tratar: la obtención de los tópicos más relevantes para tratar en el texto, que en nuestro caso son unidades de información representadas por \emph{owl:Class} y \emph{owl:NamedIndividual}; y el recorrido del grafo para obtener las unidades de información adecuada que maximicen la relación semántica respecto a los tópicos del texto. En este caso, las unidades de información que se tienen en cuenta para el recorrido del grafo son ambas \emph{owl:Class} y \emph{owl:ObjectProperty}.
 
\section{Presentación del problema de verbalización}
%Como presentamos en \ref{sec:tareas_gnl}, existen ciertas tareas a desarrollar en un sistema de generación de texto.
Como presentamos en el Objetivo, se esperan alcanzar dos cualidades principales en el texto de salida:
\begin{enumerate}
    \item A nivel macro, el texto debe tener una estructura que permita visualizar los principales temas, presentándolos de manera aislada (secciones, párrafos), pero que en conjunto se complementen para que el lector pueda establecer las relaciones adecuadas entre ellos. Como objetivo principal se espera alcanzar una estructura que presente la información más importante y necesaria cuanto antes, preparando al lector para que pueda comprender información nueva.
    \item A nivel micro, el texto debe cumplir con cierta fluidez, utilizando frases que tengan una gramática aceptable, tratando de maximizar la cohesión.
\end{enumerate}{}

\chapter{Organización de la información\label{cap:organizacion}}

Antes de comenzar con el desarrollo de los procesos asociados a la Generación de Lenguaje Natural, iniciaremos con el problema de \textit{Organizar la Información} de una ontología.  

Este capítulo describe el diseño e implementación de un módulo para organizar y estructurar la información de una ontología de forma más significativa. Reiter E.~\cite{reiter2007architecture} ha propuesto esta etapa previa de preprocesamiento para el análisis y la interpretación de datos, aunque fue aplicada en una entrada de datos en crudo (generada directamente de sensores) en lugar de una base de conocimiento. En este caso, el resultado del preprocesamiento será usado para poder abordar la coherencia global del texto.

\section{Introducción}
Considerando lo propuesto por la Teoría de Veins~\cite{cristea1998veins}, para proveer de coherencia global a un texto debemos crear una estructura jerárquica que facilite la relación entre unidades textuales. A medida que se avanza en el desarrollo de un texto, un concepto solo debería ser capaz de referenciar a otros conceptos que hayan sido nombrados anteriormente, ya sea en el mismo nivel de jerarquía o en niveles superiores. Esto implica evitar que una entidad en un nivel \emph{N} haga referencia a entidades que se encuentren en niveles \emph{N+1} (niveles inferiores). 

Suponemos que las entidades que tienen mayor impacto en la ontología son las que deberían aparecer en los niveles superiores de la jerarquía, con el objetivo de usarlas como hilo conductores para estructurar el texto. 

%Para identificar las entidades más relevantes, usaremos las medidas de centralidad de los nodos. 

A continuación analizaremos la Figura~\ref{fig:pizza.owl}, la cual contiene representada gráficamente un subconjunto de las clases (Figura~\ref{fig:jerarquilapizza}) y propiedades (Figura~\ref{fig:propiedadespizza}) de la ontología {\tt pizza.owl}\footnote{https://protege.stanford.edu/ontologies/pizza/pizza.owl}, y explicaremos cómo pretendemos estructurar su información.

Nótese que hay formas más adecuadas que otras para comenzar a transmitirle a un receptor toda la información del grafo. Algunos de los posibles tópicos que engloban o relacionan a la mayoría de los datos expuestos son ``Tipo de comidas'', o ``Ingredientes de la pizza'', siendo que, la mayor cantidad de información está relacionada a la pizza. 

Si se comenzara a estructurar el texto desde ``Helado'' o desde ``Vegetariana'', se perdería algo de coherencia global al no estar fuertemente relacionado con el resto de los tópicos.

%\begin{figure}
%\centering
%\subfigure[b]{\includegraphics[scale=0.7]{img/presentacion_problema/onto_pizza.pdf}\caption{Jerarquía de clases de la ontología pizza}}
%\subfigure[b]{\includegraphics[scale=0.7]{img/presentacion_problema/onto_pizza_properties.pdf}\caption{Jerarquía de propiedades de la ontología pizza}}
%\caption{Subconjunto de la ontología de pizzas.} \label{fig:pizza.owl}
%\end{figure}

\begin{figure}
\centering

\subfigure[Jerarquía de clases de la ontología pizza]{%\fbox{
\includegraphics[scale=0.9]{img/presentacion_problema/onto_pizza.pdf}\label{fig:jerarquilapizza}}%}

\subfigure[Jerarquía de propiedades de la ontología pizza] {\includegraphics[scale=0.9]{img/presentacion_problema/onto_pizza_properties.pdf}\label{fig:propiedadespizza}}
\caption{Subconjunto de la ontología de pizzas.} \label{fig:pizza.owl}
\end{figure}


Considerando que las ontologías de la Web Semántica son Redes Libres de Escala~\cite{zhang2008scale}, esperamos que la naturaleza de las relaciones brinden información suficiente para discriminar entre el contenido principal y el menos relevante. 

Una vez identificados los nodos sobresalientes (temas principales), se debe estructurar la información a partir de ellos. Para esto se recorre el grafo de la ontología a través de las relaciones entre los nodos.

Para este problema hemos reconocido dos sub-problemas a tratar:

\begin{itemize}
    \item La obtención de los tópicos más relevantes para tratar en el texto.%, que en nuestro caso son unidades de información representadas por \emph{owl:Class} y \emph{owl:NamedIndividual}.
    \item El recorrido del grafo para obtener las unidades de información que maximicen la relación semántica respecto a los tópicos del texto (para abordar la coherencia global).% En este caso, las unidades de información que se tienen en cuenta para el recorrido del grafo son \emph{rdfs:subClassOf}, \emph{owl:ObjectProperty} y \emph{owl:DataProperty}.
\end{itemize}

Este sub-problema de la estructuración jerárquica como Árbol de entidades serán desarrollados en las siguientes secciones.  %Nos referiremos a la estructura jerárquica como Árbol de Entidades.


\section{Diseño}
Para organizar la información, se crearon cuatro módulos que transforman la ontología de entrada en un árbol con toda la información contenida en la ontología. A continuación se describe cada módulo:
\begin{itemize}
    \item \textbf{Traductor}: recibe la entrada y la traduce a una representación interna de la ontología.
    \item \textbf{Clasificador de Entidades}: clasifica y extrae las entidades más relevantes, utilizando como criterio las medidas de centralidad.
    \item \textbf{Generador del Árbol de Entidades}: se encarga de crear los niveles del Árbol de Entidades. 
    \item \textbf{Generador de Nuevo Grupo de Entidades}: recibe una Entidad y recorre sus relaciones para crear un nuevo grupo de entidades relacionado semánticamente.
\end{itemize}

En la Figura~\ref{fig:modulos_organizador_inf} se muestra la transformación de la ontología OWL a través de los módulos. Los pasos 4 y 5 son iterativos, el proceso finaliza cuando no queden más Entidades que explorar, lo cual es definido por el Generador del Árbol de Entidades.

\begin{figure}
    \centering
    \includegraphics%[width=12cm, height=6cm]
    [scale=1]{img/organizacion_informacion/modulos_organizador_de_informacion.pdf}
    \caption{Módulos que componen el Organizador de Información}
    \label{fig:modulos_organizador_inf}
\end{figure}


\subsection{Traductor de la entrada}
El módulo Traductor crea una representación interna de la ontología cargando toda la información en una estructura que facilita el acceso y recorrido de los datos. Como algunos datos son accedidos a través de un razonador o importados de otras ontologías, se optó por cargar toda esta información una única vez para reducir el tiempo de espera cuando se requieren estos datos.

\subsection{Clasificación de las entidades más relevantes}
Una vez creada la representación de la ontología, el segundo módulo debe clasificar las entidades y extraer las más relevantes para crear el primer nivel del Árbol de Entidades. A continuación se explicará cómo son extraídas estas entidades.

\subsubsection{Criterio de clasificación}
El criterio principal para establecer una clasificación sobre las clases de una ontología, y que tal clasificación sirva para seleccionar las clases más relevantes, se basa en cuánta información posee una clase respecto al resto de las clases en la ontología. Consideramos que una clase tiene más información que otra si está presente en más cantidad de dominios de \emph{ObjectProperties}. Esto permite reconocer las clases que tengan mayor cantidad de conexiones con otras clases, y que a su vez sean el núcleo de la relación. 

Para obtener la cantidad de información de cada clase, utilizaremos el grafo subyacente a la ontología y calcularemos el \emph{outdegree}~\footnote{El \emph{outdegree} es la medida de la cantidad de conexiones salientes de un nodo.} de cada nodo que represente una clase. Para el cálculo se tendrá en cuenta únicamente la relación \emph{rdf:domain}.

Teniendo en cuenta a las clases mejor clasificadas, se puede centrar la verbalización alrededor de estas clases.

Algunas ventajas de este enfoque son:
\begin{itemize}
    \item Sencillez de implementación. Únicamente se debe recorrer el grafo calculando el valor de cada nodo. El recorrido del grafo tiene a lo sumo una complejidad polinomial.
    \item No se requiere agregar información extra al dominio.
    \item No es necesario utilizar sobre la ontología un razonador que requiera una complejidad computacional que sea intratable. La única observación es que hay que inferir el dominio y rango de las \emph{ObjectProperties}. Sin embargo, la inferencia de dominio y rango se realiza teniendo en cuenta la jerarquía de \emph{ObjectProperties}, sin necesidad de inferir clases equivalentes o disjuntas.
\end{itemize}

La eficacia de este enfoque depende fuertemente de que las \emph{ObjectProperties} más importantes tengan el dominio explícito.


\subsubsection{Seleccionando las principales clases}
\label{sec:select_class}
Una vez calculado el valor de cada clase, se pueden seleccionar las clases que superen cierto umbral para agregarlas al primer nivel de la estructura jerárquica que guiará el desarrollo del texto. En este trabajo, se recorren todas las clases y se seleccionan las que superen el valor promedio entre cantidad de propiedades y cantidad de clases.
Si ninguna supera el valor promedio, se selecciona la o las clases con el valor más alto.

Cuando se seleccionan las clases que van a representar los temas principales, puede ocurrir que se elijan clases de una misma rama en la jerarquía de clases. Esto hace que se pierda algo de semántica, pues se incluyen clases en el mismo nivel, siendo que algunas son más específicas y pueden ser alcanzadas desde sus ancestros. Para evitar este problema, se optó por eliminar las subclases que sean seleccionadas en un principio, y que tengan a una clase ancestro dentro de los temas principales. Por ejemplo, continuando con la ontología de las pizzas, si las principales clases seleccionadas son \emph{comida} y \emph{pizza}, la clase \emph{pizza} sería eliminada del conjunto, pues es subclase de \emph{comida}. 

\subsection{Generador Árbol de Entidades}
El tercer módulo se encarga de crear la estructura del Árbol, sincronizando la interacción entre los demás módulos. Su objetivo es formar la jerarquía comenzando desde las Entidades más Relevantes, y agregando nuevos niveles con ayuda del Generador de Nuevos Grupos. Cuando no quedan más relaciones que recorrer en las Entidades del último nivel, revisa que todas las clases haya sido usadas al menos una vez. Si hay clases que no fueron utilizadas, invoca un algoritmo para agregar estas clases en una rama especial llamada \emph{Otras Secciones}. Es el caso, por ejemplo, de las clases \emph{Comida} y \emph{Helado} en la ontología \emph{pizza}.

Como ejemplo, el primer y segundo nivel de la ontología de las pizzas quedaría como en la Figura~\ref{fig:macro_planning_pizza}; y el tercer nivel se muestra en la Figura~\ref{fig:macro_planning_pizza_n2}.

\begin{figure}
\centering
\begin{minipage}[c]{0.7\textwidth}
%no borrar el % de dirtree porque es necesario.
{\footnotesize 
\dirtree{%
.0 .
.1 pizza.
.2 tieneIngrediente.
.2 pizzaConCarne.
.2 pizzaConNombre.
.2 $\ldots$ las restantes subclases de pizza.
.1 comida.
.2 helado.
}}
\caption{Primer y segundo nivel del Árbol de Entidades de la ontología \emph{pizza}.}
\label{fig:macro_planning_pizza}
\end{minipage}
\end{figure}

\begin{figure}
\centering
\begin{minipage}[c]{0.7\textwidth}
%no borrar el % de dirtree porque es necesario.
{\footnotesize 
\dirtree{%
.0 .
.1 pizza.
.2 tieneIngrediente.
.3 tieneBase.
.3 tieneCobertura.
.2 pizzaConCarne.
.2 pizzaConNombre.
.3 Margherita.
.3 Napoletana.
.3 $...$ las restantes subclases de pizzaConNombre.
.2 $...$ las restantes subclases de pizza, con sus respectivas subclases.
.1 comida.
.2 helado.
}}
\caption{Tercer nivel del Árbol de Entidades de la ontología \emph{pizza}.}
\label{fig:macro_planning_pizza_n2}
\end{minipage}
\end{figure}


Se puede ver en la Figura~\ref{fig:macro_planning_pizza_n2} que las subpropiedades \emph{tieneBase} y \emph{tieneCobertura} tienen en su dominio a \emph{pizza}, la clase más cercana recorriendo sus ancestros. Si existiera por ejemplo, \emph{tieneCondimento} como una tercer subpropiedad de \emph{tieneIngrediente}, cuyo dominio no tuviera \emph{pizza}, entonces no sería listada dentro del grupo \emph{tieneIngredientes} en la rama de \emph{pizza}. Sin embargo, deben explorarse las subpropiedades de \emph{tieneCondimento}, por el caso de que tenga alguna subpropiedad que tuviera como dominio a \emph{pizza}. Por ejemplo, supongamos que \emph{tieneCondimento} tiene como subpropiedad a \emph{tieneOrégano} con dominio \emph{pizza}, en ese caso se habilita a \emph{tieneCondimento} para ser insertado en el grupo junto a  \emph{tieneCobertura} y \emph{tieneBase}. 


\subsection{Generador de Nuevos Grupos de Entidades}
\label{sec:agrupando_info}
El cuarto módulo crea nuevos grupos de entidades para agregar como nuevos niveles al Árbol. Los elementos de cada grupo serán insertados como hijos de la Entidad correspondiente.

Dada una Entidad \emph{E}, deben recorrerse las relaciones de \emph{E} en la ontología, para seleccionar Entidades que serán agregadas al nuevo grupo.
El contenido de cada grupo depende del tipo de \emph{E}.
\begin{itemize}
    \item Si \emph{E} es una Clase, el nuevo grupo estará constituido por Propiedades, Individuos y Subclases. 
    \item Si \emph{E} es una Propiedad, el nuevo grupo estará constituido por Subpropiedades cuyo dominio contenga a \emph{E}, o por las clases del Rango de \emph{E} si no posee Subpropiedades.
    \item Si \emph{E} es un Individuo, no se creará ningún grupo nuevo.
\end{itemize}

Cuando existe más de una Propiedad en un grupo, se busca reemplazarlas por propiedades en común de más alto nivel en la jerarquía de propiedades. Por ejemplo, en la ontología de las pizzas, para la clase \emph{pizza} existen dos propiedades: \emph{tieneCobertura} y \emph{tieneBase}. Como estas dos propiedades tienen a su vez una subpropiedad en común llamada \emph{tieneIngrediente}, se reemplaza a las dos subpropiedades por su ancestro \emph{tieneIngrediente}.

%En la Figura~\ref{fig:diagrama_secuencia_contentGrouping} se muestra un pseudocódigo del algoritmo usado para agrupar la información de las clases.

%\begin{figure}
%    \centering
%    \includegraphics%[width=8cm, height=7cm]
%    {img/organizacion_informacion/secuencia_contentGrouping}
%    \caption{Diagrama de secuencia para agrupar la información de las Clases}
%    \label{fig:diagrama_secuencia_contentGrouping}
%\end{figure}

\subsection{Diagrama de clases}
En la Figura~\ref{fig:diagrama_clases_organizador} se muestra el diagrama de clases correspondiente al Organizador de Información. El módulo \emph{Traductor} está representado a través de la clase \emph{OntologyManager}; el módulo Clasificador de Entidades está representado por la clase \emph{ContentClasification}; el módulo que genera el Árbol de Entidades está implementado por la clase \emph{GeneratorTreeManager}, y el módulo que crea nuevos grupos de entidades está implementado en la clase \emph{ContentGrouping}. 

Las clases \emph{OWLIntClass},  \emph{OWLClass},  \emph{OWLDataProperty},  \emph{OWLObjectProperty}, \emph{OWLIndividual}, \emph{OWLRestriction} y \emph{Ontology} son usadas para representar la estructura de la ontología. Estos objetos son creados por la clase  \emph{OntologyManager}. 

\begin{figure}
    \centering
    \includegraphics{img/organizacion_informacion/diagrama_clases_organizador_informacion.pdf}
    \caption{Diagrama de clases del Organizador de Información}
    \label{fig:diagrama_clases_organizador}
\end{figure}

\section{Implementación}
Para documentar la implementación se tomarán principalmente las cuatro clases asociadas a los módulos del Organizador de Información. Para simplificar la longitud de las clases documentadas, solo se mostrará el encabezado de los métodos con sus respectivas descripciones, y solo se incluirá la implementación de los métodos más importantes.

La estructura de la ontología se crea a través de la clase \emph{OntologyManager}. Esta clase hace uso de la librería  \emph{OWLAPI}\footnote{http://owlapi.sourceforge.net/} para cargar el archivo que contiene la ontología de entrada. El beneficio de esta librería es que también importa los archivos necesarios en caso de que hayan definiciones de tipo \emph{import} en la ontología. En el Anexo~\ref{sec:clase_ontologymanager} se muestra el método principal de \emph{OntologyManager} que contiene la llamada y la explicación de los métodos involucrados para crear la estructura de la ontología.

La clase \emph{Ontology} posee la estructura creada con \emph{OntologyManager}. También posee algunos métodos útiles para el proceso de organizar la información, los cuales se muestran en el Anexo~\ref{sec:clase_ontology}.

La clase \emph{ContentClasification} contiene los métodos necesarios para clasificar y extraer las entidades más relevantes. Su implementación se muestra en el Anexo~\ref{sec:clase_content_clasif}.

La clase \emph{ContentGrouping} posee tres métodos para generar los grupos de nuevos tópicos del Árbol de Entidades. En el Anexo~\ref{sec:clase_content_grouping} se presenta esta clase.

La clase \emph{GeneratorTreeManager} recorre recursivamente las Entidades creando las ramas del Árbol de Entidades. Se presenta un pseudocódigo en el Anexo~\ref{sec:clase_generator_tree} para explicar el recorrido.

El resto de las clases solo sirven para completar la estructura en forma de grafo de la ontología, pero no agregan métodos relevantes para el Organizador de Información.



\section{Casos de estudio}
En esta sección se presentan dos casos de estudio que muestran el comportamiento de la aplicación ante dos entradas diferentes, abarcando distintos aspectos de la solución propuesta. En el primer caso será usada la ontología de las pizzas, y en el segundo se hará uso de la ontología \emph{Wine}\footnote{https://www.w3.org/TR/owl-guide/wine.rdf}.

Al momento de realizar este trabajo no se reconoce ninguna aplicación que aborde el problema de organización de la información en una ontología, por lo que se expondrán tres casos de estudio sin la posibilidad de compararlos con otros resultados.

\subsection{Organización de Ontología Pizza}
Según la descripción de la ontología, este dominio representa las pizzas y sus coberturas, por lo que esperamos que las entidades que representen a las pizzas, a las coberturas y la información complementaria tiendan a agruparse en los niveles más altos de la jerarquía, mientras que en los niveles más bajos esperamos encontrar las clases más específicas y que tienen menor impacto sobre el entendimiento del dominio.

En el Anexo~\ref{apx:traza_pizza} se muestra una salida del programa en ejecución, en el que se ve el proceso de clasificación de las Entidades para obtener los temas principales, y luego se muestran algunas partes del proceso de creación de la estructura, donde se crean los niveles de la jerarquía de la estructura.

En la Figura~\ref{fig:caso_estudio_pizza} se puede ver el resultado de aplicar el algoritmo teniendo como entrada la ontología pizza. 

Los nombres de las entidades fueron traducidos al español, para poder generar las sentencias en lenguaje español en el siguiente capítulo, pero el idioma no afecta los resultados del algoritmo propuesto.

\begin{figure}
\begin{multicols}{2}
\begin{figure}[H]
\dirtree{%
.1 Pizza.
.2 Pizza con nombre.
.3 Margherita.
.3 Frutti di mare.
.3 Giardiniera.
.3 ....
.3 Napoletana.
.2 Pizza con carne.
.2 Pizza picante.
.2 ....
.2 Pizza vegetariana.
.2 Ingredientes.
.3 Coberturas.
.4 Cobertura de pizza.
.5 Cobertura de verduras.
.6 Cobertura de pimiento.
.6 ....
.6 Cobertura de langostinos.
.3 Bases.
.4 Base de pizza.
.5 Base gruesa.
.5 Base delgada y crujiente.
}
\end{figure}

\begin{figure}[H]
\dirtree{%
.1 Otras secciones.
.2 Domain concept.
.3 Comida.
.4 Helado.
.4 Ingrediente.
.3 Pais.
.4 America.
.4 England.
.4 Italy.
.4 France.
.4 Germany.
.2 Value partition.
.3 Picante.
.4 Poco picante.
.4 Algo picante.
.4 Muy picante.
}
\end{figure}
\end{multicols}
\caption{Resultado del Organizador de Información con la ontología \emph{Pizza}.}
\label{fig:caso_estudio_pizza}
\end{figure}

La ontología pizza cuenta con 100 clases y 8 propiedades. De las 8 propiedades, 5 fueron usadas  para clasificar las 100 clases (de las 8, 3 eran inversas a otras 3 por lo que no agregaban información nueva). El promedio de información obtenido fue 1.25, siendo la clase Pizza la única en superar este valor, lo que concuerda con el resultado esperado.

Como se puede observar en la Figura~\ref{fig:caso_estudio_pizza}, en la primer columna se encuentran los tópicos principales reconocidos por el algoritmo, y en la segunda columna se encuentra una rama especial, para aquellas clases que no habían sido utilizadas en el resto del árbol.

En la jerarquía de los tópicos principales, se puede ver como toda la información se agrupa como hijas de la entidad Pizza, describiendo los tipos de pizzas y sus ingredientes. Inmediato a los ingredientes lista las coberturas, y reconoce a las bases de pizza a la misma altura que las coberturas, asignándoles la misma importancia.

En la jerarquía de \emph{Otras Secciones}, se aprecian las demás entidades que aportan información secundaria a la descripción del dominio, que no está directamente relacionada con el dominio de las pizzas, como son los países, los tipos de picante y otras comidas.

Respecto al resultado esperado, la organización de la información es satisfactoria, ya que se asimila a la propia descripción del dominio. 

\subsection{Organización de Ontología Wine}
Esta ontología tiene como objetivo describir un dominio de vinos y comidas\footnote{\url{https://protege.stanford.edu/publications/ontology_development/ontology101-noy-mcguinness.html}}, por lo que esperamos que las entidades que se consideren más relevantes sean aquellas afines a los vinos y comidas.
En la Figura~\ref{fig:caso_estudio_wine} se puede ver el resultado de aplicar el algoritmo.

\begin{figure}
\begin{multicols}{2}
{\small
\begin{figure}[H]
\dirtree{%
.1 Wine.
.2 Italian wine.
.3 Chianti.
.4 Chianti classico.
.2 $...$ \emph{(otras subclases de Wine)}.
.2 Wines descriptor.
.3 Sugar.
.4 Wine sugar.
.5 Dry.
.5 Off dry.
.5 Sweet.
.3 Colors.
.4 Wine color.
.5 White.
.5 Rose.
.5 Red.
.3 Flavors.
.4 Wine flavor.
.5 Moderate.
.5 Strong.
.5 Delicate.
.3 Bodies.
.4 Wine body.
.5 Medium.
.5 Full.
.5 Light.
.2 From grapes.
.3 Wine grape.
.4 Chenin blanc grape.
.4 $...$ \emph{(otras subclases de Wine grape)}.
}
\end{figure}}  

\begin{figure}[H]
\dirtree{%
.1 Otras secciones.
.1 Region.
.2 Central texas region.
.2 $...$ \emph{(otras subclases de Region)}.
.1 Vintage year.
.2 Year1998.
.1 Wine descriptor.
.2 Wine taste.
.1 Non consumable thing.
.1 Vintage.
.2 Vintage years.
.1 Fruit.
.1 Winery.
.2 Chateau de meursault.
.2 $...$ \emph{(otras subclases de Winery)}.
.1 Consumable thing.
.2 Meal.
.3 Courses.
.4 Meal course.
.5 Cheese nuts dessert course.
.5 $...$ \emph{(otras subclases de Meal course)}.
.5 Foods.
.6 Edible thing.
.7 Fowl.
.8 $...$.
.7 Dessert.
.8 $...$.
.7 Meat.
.8 $...$.
.7 Seafood.
.8 $...$.
.7 Sweet fruit.
.8 Grape.
.8 $...$.
.7 Pasta.
.8 $...$.
.7 Non sweet fruit.
.5 Drinks.
.6 Potable liquid.
.7 Juice.
.7 Wine.
.8 Wines descriptor.
.8 From fruits.
}
\end{figure}

\end{multicols}
\caption{Resultado del organizador de información con la ontología Wine.}
\label{fig:caso_estudio_wine}
\end{figure}


La ontología Wine cuenta con 138 clases y 16 propiedades. De las 16 propiedades, 12 fueron usadas para verificar qué clases superan el promedio para ser seleccionadas como las principales. El promedio fue de 2.5 y la única clase que lo superó fue Wine, lo que cumple parcialmente el resultado esperado, ya que parte del objetivo de la ontología es describir el dominio de los vinos. Respecto a la sección que describe las comidas, quedó desplazada a \emph{Otras Secciones}, siendo un resultado no esperado según el objetivo de la ontología. Sin embargo, analizando manualmente la ontología, se puede apreciar que el porcentaje de  información que describe a las comidas es significativamente menor en relación a la información referida a los vinos, factor por el cual resulta aceptable que no aparezcan como sección principal.

%\subsection{Organización de Ontología SNOMED-CT}

\section{Conclusiones}
Abordamos el problema de organizar la información de una ontología, tratando de maximizar la coherencia de la estructura generada. Como presentamos en el inicio de este Capítulo, tuvimos que definir:
\begin{itemize}
    \item {\bf Las entidades que iban a representar los nodos de la estructura}: donde elegimos para el primer nivel de la jerarquía a las Clases, y para los próximos niveles a Clases, Propiedades e Individuos.
    \begin{itemize}
        \item {\bf El criterio para ponderar la importancia de las Clases para crear el primer nivel}: donde decidimos usar el criterio de centralidad de un nodo, teniendo en cuenta solo la relación \emph{rdfs:domain}.
    \end{itemize}
    \item {\bf El recorrido de la ontología para crear los siguientes niveles del Árbol}: para el cual seleccionamos las relaciones \emph{rdfs:Domain}, \emph{owl:subPropertyOf}, \emph{rdfs:Range} y \emph{rdfs:subClassOf}. El recorrido para cada Entidad visita todos sus vecinos a solo un nivel de profundidad.
\end{itemize}

Los resultados fueron satisfactorios, obteniendo estructuras que corresponden a los resultados esperados. 

%Si bien el resultado de la clasificación de las entidades mas importantes depende de la cantidad de \emph{owl:ObjectProperties} y de que sus \emph{rdfs:domain} estén correctamente establecidos, en general el hecho de que son Redes Libres de Escala es un punto a favor para este enfoque.

\chapter{Generación del documento de texto}

En este Capítulo se diseñará e implementará el proceso de generación del Documento de Texto, teniendo en mente los procesos conocidos de la GNL. Usaremos el Árbol de Entidades del Organizador de Información para dar una estructura inicial al texto, y luego verbalizaremos los Constructores de la Ontología, teniendo en cuenta la coherencia local dentro de los párrafos generados.

\section{Introducción}
Se crearon tres módulos principales teniendo como referencia las actividades de la Generación de Lenguaje Natural nombradas en la Sección~\ref{sec:tareas_gnl}. Los módulos son \emph{Macro planificación}, \emph{Micro planificación} y \emph{Realización}. Dentro de cada módulo se desarrollaron submódulos, para resolver cada problema específico. 

Elegimos organizar toda la información de la ontología generando un documento dividido en secciones, subsecciones y párrafos. 

Como presentamos en el Objetivo (Sección~\ref{Intro:objetivo}), se esperan alcanzar dos cualidades principales en el texto de salida:
\begin{enumerate}
    \item A nivel macro: el texto debe tener una estructura que permita visualizar los principales temas, presentándolos de manera aislada (secciones, párrafos), pero que en conjunto se complementen para que el lector pueda establecer las relaciones adecuadas entre ellos. Como objetivo principal se espera alcanzar una estructura que presente la información más importante y necesaria lo más temprana posible, preparando al lector para que pueda comprender información nueva.
    \item A nivel micro, el texto debe cumplir con cierta fluidez, utilizando frases que tengan una gramática aceptable, tratando de maximizar la cohesión.
\end{enumerate}

Para satisfacer las necesidades del nivel macro, usaremos como punto de partida la estructura resultante del módulo de Organización de Información, que servirá para generar un Documento Inicial. Este Documento Inicial se modificará para adoptar una estética más agradable para el lector, tratando de mantener la coherencia global de la estructura original.

Parte de la estética resulta de crear secciones y párrafos con los tópicos del texto. Dado que queremos dotar al texto con coherencia local, debemos definir qué información habrá en cada sección y párrafo, y cómo estará organizada. 

Tomando algunos postulados básicos de la Teoría de Centrado, contamos con una base inicial para elegir cómo presentar en lenguaje natural y de manera coherente la información contenida en la ontología. No es el objetivo de este trabajo implementar una solución basada estrictamente en las reglas propuestas por dicha Teoría, más bien, se pretende alcanzar una cierta coherencia local motivados por la idea que subyace en la misma.

Para mantener la coherencia local según la Teoría de Centrado, debemos evitar cambiar continuamente los focos de atención. Por este motivo, para cada tópico tratado en el texto, trataremos de capturar y presentar toda su información agrupada en párrafos consecutivos.  

Dado una entidad \emph{E}, para recuperar la información que la describe, adoptaremos dos criterios principales:
\begin{itemize}
    \item El primer criterio consiste en juntar los axiomas que describen a \emph{E} en un mismo párrafo. %Los axiomas de cada entidad están guardados en la estructura de la ontología creada en el módulo \emph{Traductor} del capítulo anterior, por lo que obtenerlos y agruparlos es trivial.
    \item El segundo criterio consiste en recorrer la estructura de grafo a partir del nodo que representa a \emph{E}, navegando por las relaciones del nodo. Este recorrido busca capturar la máxima cantidad de información posible acerca de \emph{E}. 
\end{itemize}

Estos criterios permiten desarrollar el texto tomando centralmente a cada entidad una sola vez, de manera que el foco de atención solo salte hacia otras entidades cuando no haya nada más para decir acerca de \emph{E}.
\\

En el nivel micro, debemos definir una traducción desde la gramática del lenguaje OWL hacia el lenguaje humano. Luego se modificará convenientemente la traducción para alcanzar el Documento Final, generando un texto más fluido y organizando las sentencias para maximizar la cohesión.

La relación entre el Organizador de Información y el Generador del Documento de Texto se puede ver en la Figura \ref{fig:modulos_sgln}.

\begin{figure}
    \centering
    \includegraphics{img/presentacion_problema/modulos_sgln.pdf}
    \caption{Módulos que componen el Sistema de Generación de Lenguaje Natural}
    \label{fig:modulos_sgln}
\end{figure}


Ampliando el módulo \emph{Generador de Documento de Texto} en la Figura \ref{fig:modulos_documento_final}, se puede apreciar el flujo de información entre los módulos desarrollados en este Capítulo.

\begin{figure}
    \centering
    \includegraphics{img/generacion_documento/modulos_documento_final.pdf}
    \caption{Módulos para la generación del Documento Final.}
    \label{fig:modulos_documento_final}
\end{figure}


\section{Diseño Macro Planning}
\label{sec:macro_planning}
En esta etapa se planificará la estructura del documento. Como nombramos anteriormente, la organización a priori de la información, realizada en el capítulo anterior, facilita la tarea de Macro Planificación. Para crear el documento en la etapa de Macro Planificación, se reemplazó el módulo Generador del Árbol de Entidades del Organizador de Información, por un nuevo módulo específico del Macro Planning. El motivo de este reemplazo es evitar la creación de interfaces entre el Organizador de Información y el Macro Planning, por el overhead que estas suponen, ya que la creación del Árbol de Entidades y la Macro Planificación resultan ser procesos equivalentes, por lo que pueden ser reemplazables, evitando recorrer el doble de veces la estructura jerárquica.

En la figura~\ref{fig:modulos_plan_documento} se puede ver el nuevo módulo Generador del Plan de Documento, que retorna como salida un Plan de Documento.

\begin{figure}
    \centering
    \includegraphics{img/generacion_documento/modulos_plan_documento.pdf}
    \caption{Módulos para la generación del Plan de Documento.}
    \label{fig:modulos_plan_documento}
\end{figure}

El documento de texto estará representado internamente con una estructura basada en componentes. Esta representación prepara el terreno para que, luego de la Micro Planificación, el Realizador decida el diseño final del documento.

La representación interna del documento de texto tiene en cuenta los siguientes criterios:
\begin{itemize}
    \item Cada Entidad presente en la estructura que resulta del Organizador de Información es considerada un tópico. 
    \item A cada tópico del primer nivel le corresponde una sección. Los tópicos anidados (Subclases, Individuos y Relaciones) se corresponden a subsecciones (secciones anidadas).
    \item Cada sección está compuesta de al menos un párrafo y un título.
    \item Cada párrafo trata un único tópico, el cual puede ser una clase o un individuo. También, un párrafo puede o no tener alguna sentencia. Esto se debe a que no todas las entidades tienen asociados axiomas que la describan.
\end{itemize}

\subsection{Secciones y párrafos}
Las Secciones tiene como objetivo principal delimitar la información asociada a una Entidad, segmentando el texto con el fin de mantener la coherencia global.

Dentro de una sección para una Clase, se verbalizarán sus axiomas y se enumerarán sus Individuos, Subclases y Relaciones, todo en un párrafo. Enumerar esta información (es decir, presentarla sin considerarla un nuevo tópico) resulta conveniente ya que luego se crearán subsecciones para cada uno de esos componentes.

En una sección asociada a una Relación, solo habrán subsecciones acerca de las Entidades que pertenecen a la Relación.

En una sección para Individuos, se creará un párrafo para describir sus propiedades.

\subsubsection{Subsecciones}
Las subsecciones son secciones anidadas dentro de otras secciones. Una subsección puede surgir a partir de un Individuo, Subclases o Relaciones. Cuando surge a partir de un Individuo, la subsección trata como tópico al Individuo; cuando surge a partir de una Subclase, el tópico principal es la Subclase; y cuando surge a partir de una Relación, el tópico principal resulta ser el rango de la Relación, es decir, otra Clase.

\subsection{Oraciones}
Las oraciones pertenecen a los párrafos por lo que tratan un solo tópico. Cada oración que tenga como tópico a una Clase,  aborda un solo tipo de información, por lo que existe una oración para las clases disjuntas, una oración para las clases equivalentes, y así con los demás tipos de información. Cada oración correspondientes a un Individuo contiene las propiedades con sus valores declaradas sobre el Individuo.

Al tratar un solo tipo de información por oración, se maximiza la cohesión dentro de cada oración.


\subsection{Diagrama de clases}
En la figura \ref{fig:diagrama_clases_macroplanificador} se muestra el diagrama de clases correspondiente al Macroplanificador. La clase \emph{TextManager} es la implementación del módulo Generador del Plan de Documento. Esta clase se encarga de crear la clase \emph{Text}, que es la representación interna de un documento de texto. El documento está compuesto por Secciones, Párrafos y Oraciones, cada uno de los cuales tiene asociado su respectiva clase \emph{Section}, \emph{Paragraph} y \emph{Statement}.

\begin{figure}
    \centering
    \includegraphics{img/generacion_documento/diagrama_clases_macroplanificador.pdf}
    \caption{Diagrama de clases del Macro Planning.}
    \label{fig:diagrama_clases_macroplanificador}
\end{figure}

\section{Implementación Macro Planning}
La clase principal del Macroplanificador es \emph{TextManager}. Para crear el Documento Inicial, recorre las Entidades de la misma manera que la clase \emph{GeneratorTreeManager}, explicada en la Figura~\ref{sec:clase_generator_tree}, con la diferencia que crea el objeto \emph{Text} y los objetos \emph{Section} asociados a cada Entidad. Ya que el recorrido es equivalente no se presentará su implementación, en cambio comenzaremos a mostrar las clases que componen el Documento de Texto.

La clase \emph{Text} solo contiene los objetos del primer nivel del Documento de Texto. Tiene acceso a las secciones principales y a las secciones que contienen la información que no había sido utilizada en otras secciones. 

La clase \emph{Section} contiene la información referida a una Sección, como título, párrafos, sus subsecciones y el tópico asociado. Implementa el método de realización del Documento Final, el cual se verá en el Capítulo de Realización.

Cada \emph{Section} se encarga de crear el \emph{Paragraph} correspondiente a su Tópico.

La clase \emph{Paragraph} contiene una lista de todas las \emph{Statement} referidas a su Tópico. 

Cada \emph{Paragraph} se encarga de crear las \emph{Statement} correspondientes para el Tópico que recibe, y asociarles los componentes (Axiomas) que describen al Tópico. Adicionalmente el párrafo recibe el lenguaje del texto, por lo que debe instanciar las \emph{Statement} adecuadas, ya sea \emph{StatementEnglish} o \emph{StatementSpanish}.

La clase \emph{Statement} contiene la información que tendrá cada oración. Será usada por el Microplanificador para verbalizar la información e implementar los métodos de la etapa de Micro Planning. 

\section{Diseño Micro Planning}
En esta estapa se llevarán a cabo las tareas para describir una Entidad en lenguaje natural a partir de la información que contiene cada \emph{Statement} creada en el Macro Planning. Esta descripción en lenguaje natural se alcanza a través de una o varias oraciones producidas con una estructura que se adapte a la gramática del lenguaje español.

Como existen diversas estructuras sintácticas para expresar un mismo contenido, a modo de simplificación se asoció a cada constructor OWL solo algunas formas de verbalización. Se trató de generar oraciones que sean fieles al significado expresado en los axiomas y de evitar ambigüedades.

Para poder trabajar con patrones de la gramática del lenguaje humano (ya sea español o inglés), se utilizó un \emph{POS Tagger}\footnote{Un \emph{Part-Of-Speech Tagger} es un software para asignar la categoría gramatical a las palabras de un texto.} para etiquetar las palabras.

\subsection{Verbalización de una oración}
Tanto las oraciones acerca de Clases como las de Individuos requieren la verbalización de los constructores OWL. Solo en el caso de las oraciones acerca de Clases requieren además la enumeración de sus instancias, propiedades y subclases. Con enumeración de las propiedades (y de instancias y subclases por igual), nos referimos únicamente a presentarlas, por ejemplo, las propiedades \emph{tieneCobertura} y \emph{tieneBase} para la clase Pizza resultan en la oración ``una pizza tiene base de pizza y tiene cobertura de pizza''.

Las tareas necesarias para la verbalización de una oración son:
\begin{itemize}
    \item Verbalizar los constructores OWL o enumerar las propiedades, instancias y subclases.
    \item Eliminar información redundante entre oraciones a través del proceso de agragación.
    \item En el caso de las Clases, se reemplaza con expresiones de referencia a los nombres de las Clases en oraciones consecutivas evitando que se repitan y que genere un texto poco fluido.
\end{itemize}

\subsection{Sintaxis OWL 2}
\label{sec:gen_doc_sintaxis_owl}
Teniendo en cuenta la estructura jerárquica de la sintaxis (a partir de su gramática BNF y de los diagramas de clases de su documentación~\cite{OWL2W3C}), agrupamos en diferentes niveles de abstracción las tres categorías sintácticas: en el nivel inferior se encuentran las Entidades, en el nivel intermedio las Expresiones de Clases, y en el nivel superior los Axiomas. Esta separación en niveles de abstracción será útil para organizar la verbalización. El conjunto de constructores elegido para verbalizar se encuentra en las siguientes listas: en la lista~\ref{list:constructores_axiomas} se agrupan los constructores del nivel de Axiomas, en la lista~\ref{list:constructores_expresiones} los constructores de Expresiones de Clases y en la lista~\ref{list:constructores_entity} los constructores de nivel de Entidades.
\begin{figure}
\begin{multicols}{2}
\captionof{listCap}{Constructores de Entidades}
\label{list:constructores_entity}
    \begin{itemize}
        \item owl:class
        \item owl:objectProperty
        \item owl:dataProperty
        \item owl:individual
        \item[\vspace{\fill}]
    \end{itemize}

\captionof{listCap}{Constructores de Axiomas}
\label{list:constructores_axiomas}
    \begin{itemize}
        \item rdfs:subClassOf
        \item owl:equivalentClass
        \item owl:disjointWith
        \item rdfs:domain
        \item rdfs:range
    \end{itemize}
    \end{multicols}
\end{figure}


\begin{figure}
\captionof{listCap}{Constructores de Expresiones de Clase}
\label{list:constructores_expresiones}
    \begin{itemize}
    \begin{multicols}{2}
        \item owl:intersectionOf
        \item owl:unionOf
        \item owl:complementOf
        \item owl:allValuesFrom
        \item owl:someValuesFrom 
        \item owl:minCardinality
        \item owl:maxCardinality
        \item owl:Cardinality
        \item owl:oneOf
        \item owl:hasValue
        \end{multicols}
    \end{itemize}
\end{figure}

\subsection{De OWL 2 a lenguaje natural}
La verbalización de la sintaxis OWL 2 se corresponde a la tarea de traducir el lenguaje OWL al lenguaje humano. Para llevar a cabo la traducción, se tendrán en cuenta los niveles de abstracción presente en la Sección~\ref{sec:gen_doc_sintaxis_owl}. 

El enfoque adoptado para crear la sintaxis de las oraciones se basa en un recorrido bottom-up de la jerarquía de un Axioma. Se comienza desde las hojas, construyendo oraciones parciales a partir de las Entidades, luego se procede a subir por los niveles a través de los constructores de Expresiones de Clases, componiendo nuevas oraciones parciales, hasta alcanzar la raíz de la jerarquía, donde se termina de construir la oración final. Adicionalmente, se realizan algunos tratamientos morfológicos para agregar fluidez y coherencia al texto, tal como insertar artículos, y hacer concordar los sustantivos en género y número.

Con el objetivo de comenzar a formar las oraciones parciales, el primer paso a realizar es verbalizar las Entidades, ya que están directamente relacionadas con las IRIs (o tienen acceso al \emph{label}, en caso de extraer los nombres desde los \emph{labels}).

En la Gramática~\ref{gram:entity} se muestra una porción de la gramática BNF del lenguaje OWL 2 asociada a las Entidades.
\begin{GrammarEnv}
%detecta un error de sintaxis pero es mentira.
\begin{grammar}
[(colon){$\rightarrow$}]
[(semicolon)$|$]
[(comma){}]
[(period){\vspace{0.3cm} \\}]
[(quote){\begin{bf}}{\end{bf}}]
[(nonterminal){$<$}{$>$}]
%<expression> : <number> ; <number>, [\{"asd"\}], <relational\_operator>, <number>.
%<number> : <digit> ; <digit> , <number>.
%<digit> : "0";"1";"2";"3";"4";"5";"6";"7";"8";"9".
%<relational\_operator> : $"="$;"$\lessthan \greaterthan$";"$\lessthan$";"$\greaterthan$"; "$\lessthan=$";"$\greaterthan=$";"in".
\fbox{\begin{minipage}{14cm} \small{ 
<Entity> : <Class> ; <ObjectProperty> ; <DataProperty> ; <NamedIndividual> ; <AnnotationProperty>.
<Class> : <IRI> .
<ObjectProperty> : <IRI> .
<DataProperty> : <IRI> .
<AnnotationProperty> : <IRI> .
<NamedIndividual> : <IRI> .
}
\end{minipage}
}
\end{grammar}
\caption{Porción de gramática asociada a las Entidades.}\label{gram:entity}
\end{GrammarEnv}

En el siguiente nivel se tienen los constructores de Expresiones de Clase, con las cuales es posible generar oraciones complejas subordinando las expresiones que se encuentran anidadas, o componer nuevas oraciones parciales teniendo en cuenta los tipos de componentes involucrados.

En la Gramática~\ref{gram:expr_clases} se muestra a modo de ejemplo una sección de la gramática con constructores de este nivel.
\begin{GrammarEnv}
\begin{grammar}
[(colon){$\rightarrow$}]
[(semicolon)$|$]
[(comma){}]
[(period){\vspace{0.3cm} \\}]
[(quote){\begin{bf}}{\end{bf}}]
[(nonterminal){$<$}{$>$}]
%<expression> : <number> ; <number>, [\{"asd"\}], <relational\_operator>, <number>.
%<number> : <digit> ; <digit> , <number>.
%<digit> : "0";"1";"2";"3";"4";"5";"6";"7";"8";"9".
%<relational\_operator> : $"="$;"$\lessthan \greaterthan$";"$\lessthan$";"$\greaterthan$"; "$\lessthan=$";"$\greaterthan=$";"in".
\fbox{\begin{minipage}{14cm} \small{ 
<ClassExpression> : <Class> ; <ObjectIntersectionOf> ; <ObjectSomeValuesFrom> .
<ObjectSomeValuesFrom> : "ObjectSomeValuesFrom" "(" <ObjectPropertyExpression> <ClassExpression> ")" .
<ObjectPropertyExpression> : <ObjectProperty> .
}
\end{minipage}}
\end{grammar}
\caption{Porción de gramática asociada a las Expresiones de Clases.}\label{gram:expr_clases}
\end{GrammarEnv}

Puede apreciarse en la gramática que existe recursividad entre los constructores, lo que produce que un constructor aparezca en diferentes niveles y pueda combinarse con cualquier otro constructor del nivel de Expresiones de Clases. En este punto es importante que cada constructor pueda componerse con los demás, para asegurar cualquier combinación posible.

Por último en el nivel más abstracto se encuentran los constructores de Axiomas. Estos constructores tienen la característica de no ser recursivos entre ellos, por lo que es posible generar oraciones independientes, yuxtapuestas o coordinadas. Un ejemplo de estos constructores se ve en la Gramática~\ref{gram:axiom}.

\begin{GrammarEnv}
\begin{grammar}
[(colon){$\rightarrow$}]
[(semicolon)$|$]
[(comma){}]
[(period){\vspace{0.3cm} \\}]
[(quote){\begin{bf}}{\end{bf}}]
[(nonterminal){$<$}{$>$}]
%<expression> : <number> ; <number>, [\{"asd"\}], <relational\_operator>, <number>.
%<number> : <digit> ; <digit> , <number>.
%<digit> : "0";"1";"2";"3";"4";"5";"6";"7";"8";"9".
%<relational\_operator> : $"="$;"$\lessthan \greaterthan$";"$\lessthan$";"$\greaterthan$"; "$\lessthan=$";"$\greaterthan=$";"in".
\fbox{\begin{minipage}{14cm} \small{ 
<ClassAxiom> : <SubClassOf> ; <EquivalentClasses> ; <DisjointClasses> ; <DisjointUnion> .
<EquivalentClasses> : "EquivalentClasses" "(" <axiomAnnotations> <ClassExpression> <ClassExpression> \{ <ClassExpression> \} ")".
}
\end{minipage}}
\end{grammar}
\caption{Porción de gramática asociada a los Axiomas.}\label{gram:axiom}
\end{GrammarEnv}

\subsection{Componentes y oraciones parciales}
Durante la creación de una oración, se recorren los constructores de Expresiones de Clases y de Entidades, creando oraciones parciales y componiéndolas entre ellas. Dado que cada constructor puede componerse con cualquier otro y en cualquier nivel de profundidad, se decidió asociar a cada constructor un tipo de componente oracional, de esta manera se evita una exhaustiva programación de composiciones.
Los tipos de componentes son los siguientes:
\begin{itemize}
    \item Término (T): caracterizado por no poseer verbo. Pueden contener adverbios, sustantivos y adjetivos.
    \item Sintagma Verbal (SV): debe poseer un verbo.
    \item Oración Negativa (ON): representa la negación de un componente.
    \item Unión: representa una disyunción de oraciones parciales.
    \item Intersección: representa una adición o subordinación de oraciones parciales.
\end{itemize}

Estos tipos de componentes definen un sistema de tipos, en el que cada tipo puede componerse con otro y dan como resultado un nuevo componente. El tipo de una composición depende del constructor que opere con esos componentes.
%En la mayoría de los casos el tipo de componente es predecible, a excepción de la \emph{intersección}.

Los componentes que retorna cada constructor se definen a continuación:
\begin{itemize}
    \item Componente Término (T):
    \begin{itemize}
        \item Clase.
        \item Individuo
    \end{itemize}
    \item Componente Sintagma Verbal (SV):
    \begin{itemize}
        \item Propiedad.
        \item Constructores de cuantificación.
        \item Constructor \emph{hasValue}.
        \item Constructores con cardinalidad.
    \end{itemize}
    \item Componente Union: 
    \begin{itemize}
        \item \emph{UnionOf}.
        \item \emph{OneOf}
    \end{itemize}
    
    \item Componente Intersección: solo el constructor \emph{IntersectionOf}.
    \item Componente Oración Negativa (ON): solo el constructor \emph{ComplementOf}.
\end{itemize}

El beneficio de este sistema es que cada constructor sabe cómo realizar la composición, en función del tipo de cada uno de sus operandos, lo que reduce la cantidad de combinaciones a programar (en contraste con tener que programar el proceso de composición de cada constructor teniendo en cuenta que sus operandos serían otros constructores en lugar de componentes). 

Además, incorporar un sistema como este, evita el uso de otras técnicas menos afines a la lingüística, como \emph{templates} prediseñados o texto genérico para cada caso particular, ya que se basa en el uso de patrones más cercanos a la Lingüística Teórica, apoyando el uso de teorías lingüísticas en el campo del Procesamiento de Lenguaje Natural.


Sin embargo, estos componentes no intentan ser exhaustivos ni completamente fieles a la sintaxis del lenguaje humano, sino que intentan acaparar de la forma más general los posibles resultados de cada constructor, para mantener una sintaxis sencilla y aceptable. Una Clase, por ejemplo, podría cumplir la función de sintagma nominal, adjetival o adverbial, pero para simplificar la nomenclatura, se retorna un tipo más genérico y al momento de usar una Clase, de ser necesario, se resuelve la composición en función de los elementos que la constituyen. Como veremos más adelante, a veces es innecesario discriminar los tipos de sintagmas.


\subsection{Convención y suposiciones del nombrado de Entidades}
Para no limitar la aplicación a una convención de nombrado, no se requiere una forma estricta para nombrar las entidades en una ontología. Sin embargo, nos basamos en algunas suposiciones que resultan natural al crear una ontología:
\begin{enumerate}
    \item Las entidades pueden tener su nombre ya sea en la IRI o en un label\footnote{Estas opciones son excluyentes, al verbalizar una ontología solo se extraen los nombres de un solo lugar, por lo que todos los nombres deben ser extraídos de IRI o de label sin combinarse}. Si el nombre es extraído de una IRI, debe estar escrito con el estilo CamelCase. Si el nombre se extrae de un label, debe estar escrito en lenguaje natural, y debe tener el idioma del label (en este caso, español ``es'' o ingles ``en'').
    \item No hay condiciones para el nombre de una clase, sin embargo es preferible que contenga al menos un sustantivo.
    \item No hay condiciones para el nombre de una propiedad, sin embargo es preferible que contenga al menos un verbo. 
    
    No se infieren verbos, por lo que si se quiere expresar, por ejemplo, que una clase \emph{está localizada en} algún lugar, la propiedad debería llamarse \emph{estáLocalizadaEn} y no \emph{localizadaEn}. Al igual que los verbos, no es conveniente omitir las preposiciones finales, por ejemplo, en la propiedad \emph{esParte}, es recomendable utilizar \emph{esParteDe}.
\end{enumerate}

Estas condiciones ayudan a mejorar la fluidez del texto y no suponen una gran carga en los desarrolladores de ontologías.

De hecho, existen varios trabajos que abordan la nomenclatura de conceptos, especificando convenciones de nombrado para el desarrollo de las ontologías, coincidiendo en algunos aspectos (y siendo hasta más estrictos) con nuestros requerimientos~\cite{montiel2011style}~\cite{OntoCheck}.

\subsection{Tratamiento sintáctico de los nombres de Entidades}
Para facilitar el entendimiento de las secciones referidas al micro planning, se explicará cómo fueron tratados los nombres de las entidades, desde el punto de vista sintáctico. 

Dado que los nombres de las Entidades pueden no corresponderse fielmente con la gramática de un lenguaje (por falta de palabras funcionales, tales como artículos o preposiciones), se optó por no realizar un análisis sintáctico a los nombres. Sin embargo, para soportar la creación de patrones gramaticales, se etiquetaron las palabras de cada Entidad para reconocer la función gramatical de cada una y poder llevar a cabo las composiciones.

Se consideró, a cada conjunto de palabras que conforman el nombre de una entidad, como un componente incompleto (oración parcial) que es parte de una oración más grande. Por este motivo, se divide en dos partes fundamentales: una parte inicial y un complemento, entre los cuales pueden insertarse nexos y cuantificaciones. A continuación se explican cómo se reconocen ambas partes en las clases y propiedades:
\begin{itemize}
    \item Los nombres de las propiedades (en los que suponemos existe al menos un verbo), se considera como parte inicial a todas las palabras desde el inicio hasta el primer verbo, incluído el verbo. El complemento se conforma con todas las palabras que siguen al verbo. Por ejemplo: la propiedad \emph{tiene color} se compone de parte inicial: \emph{tiene}, y complemento: \emph{color}. La propiedad \emph{se solapa con} se compone de \emph{se solapa} como parte de inicial, y \emph{con} como complemento.
    \item Los nombres de las clases (en los que suponemos existe al menos un sustantivo), se considera como parte inicial a todas las palabras desde el inicio hasta el primer sustantivo, incluido el sustantivo. El complemento se conforma con todas las palabras que siguen al sustantivo. Por ejemplo: la clase \emph{agente social}, tiene como parte inicial \emph{agente} y complemento \emph{social}. La clase \emph{cobertura de pizza} se compone de \emph{cobertura} con complemento \emph{de pizza}.
    \item Los nombres de los individuos son tratados igual que los nombres de las clases.
\end{itemize}


\subsection{Verbalización de constructores OWL}
\label{sec:verbalizacion_constructores}
Esta tarea se encarga de componer las oraciones según la información de cada constructor OWL.
%Las oraciones pueden enfocarse ya sea en una \emph{owl:Class} o  \emph{owl:NamedIndividual}. Cuando se enfocan en una \emph{owl:Class}, pueden contener la siguiente información: clases equivalentes, disjuntas y superclases. Cuando se enfoca en un \emph{owl:NamedIndividual}, puede contener información acerca de los valores de sus propiedades.
%Para cada uno de estos tipos de información, se pueden presentar una o varias oraciones. Las oraciones que tratan el mismo tipo de información se organizan adyacentemente entre ellas.
En algunos constructores, tener un único patrón gramatical resulta insuficiente, debido a que pueden recibir distintas estructuras gramaticales como entrada, que requieren diferentes formas de componerse. Por lo tanto, en estos casos se diseñó más de una forma de composición, mejorando la variabilidad, fluidez e interpretación de las oraciones. 

Para reconocer cómo componer las oraciones en cada constructor, se revisó empíricamente los axiomas de algunas ontologías, y se buscó utilizar oraciones que sean lo más genéricas posibles, que permitan la comprensión de los axiomas. 

A continuación se explican los patrones gramaticales usados para cada constructor. %El signo $+$ corresponde a la concatenación de cadenas de texto.

\subsubsection{Restricción de cardinalidad \emph{ObjectCardinalityRestriction}}

La gramática~\ref{gram:object_card_rest} muestra los patrones diseñados para las restricciones de cardinalidad sobre \emph{objectProperty}. En los casos donde sea posible, si la clase sobre la que se aplica la restricción es \emph{owl:Thing}, se reemplaza por el rango de la propiedad, el complemento de la propiedad (como el patron$_2$) o por último por la palabra ``cosa''.

El enlace es seleccionado dependiendo del operador \emph{MinCardinality}, \emph{MaxCardinality} y \emph{ExactCardinality}.

Ejemplo: para la clase ``pizza interesante'', se tiene el axioma: ``tieneCobertura min 3 owl:Thing'', que puede verbalizarse como ``tiene como mínimo 3 coberturas de pizza''. La parte de la oración que se corresponde a \emph{coberturas de pizza}, es extraída del rango de la propiedad.

\begin{GrammarEnv}
\begin{grammar}
[(colon){$\rightarrow$}]
[(semicolon)$|$]
[(comma){}]
[(period){\vspace{0.3cm} \\}]
[(quote){\begin{bf}}{\end{bf}}]
[(nonterminal){$<$}{$>$}]
\fbox{\begin{minipage}{14cm} \small{ 
<RestriccionCard> : <patron$_1$> ; <patron$_2$> ; <patron$_3$> ; <patron$_4$> .
<patron$_1$> : <parteInicialPropiedad> <complementoPropiedad> <enlace> <cardinalidad> <clases> .
<patron$_2$> : <parteInicialPropiedad> <enlace> <cardinalidad> <complementoPropiedad>.
<patron$_3$> : <parteInicialPropiedad> <enlace> <cardinalidad> <rangoPropiedad>.
<patron$_4$> : <parteInicialPropiedad> <enlace> <cardinalidad> <clases>.
<enlace> : "como mínimo" ; "como máximo"; "exactamente".
}
\end{minipage}}
\end{grammar}
\caption{Patrones para ObjectCardinalityRestriction.}\label{gram:object_card_rest}
\end{GrammarEnv}

\subsubsection{Restricción de cardinalidad \emph{DataCardinalityRestriction}}
La gramática~\ref{gram:data_card_rest} muestra los patrones diseñados para las restricciones de cardinalidad sobre \emph{dataProperty}. El patron$_2$ es específico para cuando la propiedad no tiene verbo. 

\begin{GrammarEnv}
\begin{grammar}
[(colon){$\rightarrow$}]
[(semicolon)$|$]
[(comma){}]
[(period){\vspace{0.3cm} \\}]
[(quote){\begin{bf}}{\end{bf}}]
[(nonterminal){$<$}{$>$}]
\fbox{\begin{minipage}{14cm} \small{ 
<RestriccionCard> : <patron$_1$> ; <patron$_2$>.
<patron$_1$> : <parteInicialPropiedad> <enlace> <cardinalidad> <complementoPropiedad> .
<patron$_2$> : "tiene" <enlace> <cardinalidad> <nombrePropiedad>.
<enlace> : "al menos" ; "como máximo"; "exactamente".
}
\end{minipage}}
\end{grammar}
\caption{Patrones para DataCardinalityRestriction.}\label{gram:data_card_rest}
\end{GrammarEnv}

\subsubsection{Restricción de cuantificación \emph{QuantifiedObjectRestriction}}
La gramática~\ref{gram:quant_obj_rest} muestra los patrones diseñados para las cuantificaciones sobre \emph{objectProperty}.

\begin{GrammarEnv}
\begin{grammar}
[(colon){$\rightarrow$}]
[(semicolon)$|$]
[(comma){}]
[(period){\vspace{0.3cm} \\}]
[(quote){\begin{bf}}{\end{bf}}]
[(nonterminal){$<$}{$>$}]
\fbox{\begin{minipage}{14cm} \small{ 
<RestriccionQuant> : <patron$_1$> ; <patron$_2$> ; <patron$_3$> ; <patron$_4$> ; <patron$_5$> .
<patron$_1$> : <verboPropiedad> "como" <complementoPropiedad> <enlace> <clases> .
<patron$_2$> : <verboPropiedad> <enlace> <clases> .
<patron$_3$> : <verboPropiedad> <enlace> "a" <clases> .
<patron$_4$> : <verboPropiedad> "a" <enlace> <clases> .
<patron$_5$> : <verboPropiedad> <complementoPropiedad> <enlace> <clases> .
<enlace> : "algún" ; "alguna"; "algo"; "exclusivamente" .
}
\end{minipage}}
\end{grammar}
\caption{Patrones para QuantifiedObjectRestriction.}\label{gram:quant_obj_rest}
\end{GrammarEnv}

Los enlaces dependen del tipo de cuantificador y del tipo de palabra que lo proceda. Si el cuantificador es existencial, el enlace sería ``alguna/algún'' para palabras que sean sustantivo femenino o masculino respectivamente, o con cualquier otra palabra sería ``algo''.  Para el cuantificador universal el enlace es siempre ``exclusivamente''.

Ejemplo: la clase ``Margherita'' que es subclase de ``PizzaConNombre'', tiene el axioma ``tieneCobertura only 
    (CoberturaDeMozzarella or CoberturaDeTomate)'' (siendo only el cuantificador universal), el cual se traduce a ``tiene cobertura de mozzarella o tomate''. También posee dos axiomas equivalentes al anterior: ``tieneCobertura some CoberturaDeMozzarella'' y ``tieneCobertura some CoberturaDeTomate'', (siendo some el cuantificador existencial), los cuales se traducen a ``tiene alguna cobertura de mozzarella'' y ``tiene alguna cobertura de tomate''.


\subsubsection{Restricción de cuantificación \emph{QuantifiedDataRestriction}}
La gramática~\ref{gram:quant_data_rest} muestra los patrones diseñados para las cuantificaciones sobre \emph{dataProperty}.

\begin{GrammarEnv}
\begin{grammar}
[(colon){$\rightarrow$}]
[(semicolon)$|$]
[(comma){}]
[(period){\vspace{0.3cm} \\}]
[(quote){\begin{bf}}{\end{bf}}]
[(nonterminal){$<$}{$>$}]
\fbox{\begin{minipage}{14cm} \small{ 
<RestriccionQuant> : <patron$_1$> ; <patron$_2$> .
<patron$_1$> : <verboPropiedad> <enlace> <complementoPropiedad> "de tipo" <rangoPropiedad> .
<patron$_2$> : <verboPropiedad> <complementoPropiedad> <enlace> "de tipo" <rangoPropiedad> .
<enlace> : "algún" ; "alguna"; "algo"; "exclusivamente" .
}
\end{minipage}}
\end{grammar}
\caption{Patrones para QuantifiedDataRestriction.}\label{gram:quant_data_rest}
\end{GrammarEnv}

\subsubsection{Restricción \emph{OWLObjectComplementOf}}
La gramática~\ref{gram:complement_rest} muestra los patrones diseñados para el complemento de una Expresión de Clase.

\begin{GrammarEnv}
\begin{grammar}
[(colon){$\rightarrow$}]
[(semicolon)$|$]
[(comma){}]
[(period){\vspace{0.3cm} \\}]
[(quote){\begin{bf}}{\end{bf}}]
[(nonterminal){$<$}{$>$}]
\fbox{\begin{minipage}{14cm} \small{ 
<RestriccionComplement> : <patron$_1$>.
<patron$_1$> : <enlace> <clase> .
<enlace> : "lo opuesto de" ; "no"; "ni"; "excepto" .
}
\end{minipage}}
\end{grammar}
\caption{Patrones para OWLObjectComplementOf.}\label{gram:complement_rest}
\end{GrammarEnv}

\subsubsection{Restricción \emph{OWLObjectInverseOf}}
La gramática~\ref{gram:inverse_rest} muestra los patrones diseñados para el inverso de una Propiedad.

\begin{GrammarEnv}
\begin{grammar}
[(colon){$\rightarrow$}]
[(semicolon)$|$]
[(comma){}]
[(period){\vspace{0.3cm} \\}]
[(quote){\begin{bf}}{\end{bf}}]
[(nonterminal){$<$}{$>$}]
\fbox{\begin{minipage}{14cm} \small{ 
<RestriccionInverse> : <patron$_1$>.
<patron$_1$> : <enlace> <Propiedad> .
<enlace> : "es lo opuesto de" .
}
\end{minipage}}
\end{grammar}
\caption{Patrones para OWLObjectInverseOf.}\label{gram:inverse_rest}
\end{GrammarEnv}

\subsubsection{Restricción \emph{OWLObjectHasValue}}
La gramática~\ref{gram:hasvalue_rest} muestra los patrones diseñados para los valores de una Propiedad.

\begin{GrammarEnv}
\begin{grammar}
[(colon){$\rightarrow$}]
[(semicolon)$|$]
[(comma){}]
[(period){\vspace{0.3cm} \\}]
[(quote){\begin{bf}}{\end{bf}}]
[(nonterminal){$<$}{$>$}]
\fbox{\begin{minipage}{14cm} \small{ 
<RestriccionHasValue> : <patron$_1$> ; <patron$_2$>.
<patron$_1$> : <parteInicialPropiedad> <complementoPropiedad> <Individuo>.
<patron$_2$> : <parteInicialPropiedad> <Individuo> .
}
\end{minipage}}
\end{grammar}
\caption{Patrones para OWLHasValue.}\label{gram:hasvalue_rest}
\end{GrammarEnv}

\subsubsection{Restricción \emph{OWLObjectHasSelf}}
La gramática~\ref{gram:hasself_rest} muestra los patrones diseñados para una auto restricción de una Propiedad sobre sí misma.

\begin{GrammarEnv}
\begin{grammar}
[(colon){$\rightarrow$}]
[(semicolon)$|$]
[(comma){}]
[(period){\vspace{0.3cm} \\}]
[(quote){\begin{bf}}{\end{bf}}]
[(nonterminal){$<$}{$>$}]
\fbox{\begin{minipage}{14cm} \small{ 
<RestriccionHasSelf> : <patron$_1$>.
<patron$_1$> : <parteInicialPropiedad> <complementoPropiedad> "Sí MISMO".
}
\end{minipage}}
\end{grammar}
\caption{Patrones para OWLObjectHasSelf.}\label{gram:hasself_rest}
\end{GrammarEnv}


\subsubsection{Restricción N-Ary \emph{UnionOf} y \emph{oneOf}} 
Como representan una secuencia de elementos, en la que uno o más elementos pueden ser verdaderos, se eligió concatenar a todos los elementos con el conector disyuntivo ``o''. Sean e1, e2,...,eN los elementos de la enumeración, el patrón gramatical elegido es: ``e1, e2, ..., eN-1 \emph{o} eN''.

\subsubsection{Restricción N-Ary \emph{InterseccionOf}}
La intersección da lugar tanto a la adición como a la subordinación de oraciones. Antes de verbalizar una intersección, se ordenan sus componentes de mayor a menor prioridad, teniendo en cuenta que: T, ON, Unión, Intersección tienen la misma prioridad, y SV tiene menor prioridad. 

Para determinar cómo verbalizar una intersección, se tienen en cuenta las siguientes condiciones:
\begin{itemize}
    \item Si la intersección tiene dos operandos:
        \begin{itemize}
            \item Si el primero es de tipo T  y el segundo SV, se subordina la segunda oración de la siguiente manera:
            ``contenido primer operando + ``que'' + contenido segundo operando''
            \item en cualquier otro caso, se coordinan los operandos:
            ``contenido primer operando + ``y'' + contenido segundo operando''
        \end{itemize}
        \item Si la intersección tiene más de dos operandos, se agrupan según el tipo  (T, SV, ON) y luego se coordinan con la conjunción ``y''.
\end{itemize}

Ejemplo, el axioma de clase equivalente ``Pizza and (tieneCobertura some (CoberturaDePizza and (tienePicante some MuyPicante)))'' de la clase pizzaPicanteEquivalente, contiene dos intersecciones las cuales poseen los mismos componentes: T$+$SV, por lo que utilizamos el conector ``que'' para subordinar el segundo operando. Con este criterio, el axioma anterior sería ``pizza picante es una pizza \emph{que} tiene alguna cobertura de pizza \emph{que} tiene algún picante muy picante''.

En caso de asociar la intersección al conector \emph{y} sin tener en cuenta el contexto, se generaría una oración como la siguiente ``pizza picante es una pizza \emph{y} tiene alguna cobertura de pizza \emph{y} tiene algo muy picante'', dando lugar a ambigüedad, pues la segunda ``y'' podría cumplir la función de adición, interpretándose que la pizza tiene algo picante, en lugar de cumplir una función de subordinación, donde la interpretación correcta es que la cobertura de la pizza es picante. 

\subsection{Enumeración de propiedades, subclases e individuos}
Enumerar esta información es solo para anunciar la existencia de estos componentes, por lo que la sintaxis es simple.

Para las subclases se creó el patrón de la Gramática~\ref{gram:subclases}.
Para los individuos se creó un patrón similar al de las subclases.
Para las propiedades se creó el patrón de la Gramática~\ref{gram:propiedades}.

\begin{GrammarEnv}
\begin{grammar}
[(colon){$\rightarrow$}]
[(semicolon)$|$]
[(comma){}]
[(period){\vspace{0.3cm} \\}]
[(quote){\begin{bf}}{\end{bf}}]
[(nonterminal){$<$}{$>$}]
\fbox{\begin{minipage}{14cm} \small{ 
<EnumerarSubclases> : <patron$_1$> ; <patron$_2$>.
<patron$_1$> : "existen las siguientes clases de" <topico> "$\colon$" <subclases>.
<patron$_2$> : <subclase> "es la única clase de" <topico>.
}
\end{minipage}}
\end{grammar}
\caption{Patrones para enumerar subclases.}\label{gram:subclases}
\end{GrammarEnv}

\begin{GrammarEnv}
\begin{grammar}
[(colon){$\rightarrow$}]
[(semicolon)$|$]
[(comma){}]
[(period){\vspace{0.3cm} \\}]
[(quote){\begin{bf}}{\end{bf}}]
[(nonterminal){$<$}{$>$}]
\fbox{\begin{minipage}{14cm} \small{ 
<EnumerarPropiedades> : <patron$_1$> ; <patron$_2$> ; <patron$_3$>.
<patron$_1$> : <parteInicialPropiedad> <complementoPropiedad> <rangoPropiedad> .
<patron$_2$> : <parteInicialPropiedad> "como" <complementoPropiedad> "a" <rangoPropiedad> .
<patron$_3$> : <parteInicialPropiedad> <rangoPropiedad> .
}
\end{minipage}}
\end{grammar}
\caption{Patrones para enumerar propiedades.}\label{gram:propiedades}
\end{GrammarEnv}



\subsection{Expresiones de referencia}
El proceso de referenciación se ve facilitado por el uso de la Teoría de Centrado. 
Teniendo en cuenta que cada párrafo (y por lo tanto cada oración que lo compone) trata centralmente una única Clase o Individuo, no se intercalan tópicos como focos de atención, por lo que se reduce la posibilidad de ambigüedad. Las expresiones de referencia solo fueron implementadas para las Clases.

Para llevar a cabo la referenciación, se tuvo en cuenta el uso de pronombres demostrativos (``este/esta''), a veces, acompañados por el primer sustantivo (si posee) del nombre de la Clase. Por ejemplo: sean las oraciones ``Una Rosa es una pizza con nombre. Una Rosa tiene como cobertura a alguna cobertura de gorgonzola, mozzarella Y tomate.'' en la segunda oración en lugar de repetir ``una Rosa'', se reemplaza por una expresión de referencia ``Una Rosa es una pizza con nombre. \emph{esta} tiene como cobertura a alguna cobertura de gorgonzola, mozzarella Y tomate.''. Otro ejemplo en el que el pronombre es acompañado por el sustantivo de la Clase es: ``una cobertura de salsa de chile tabasco es una cobertura de salsa. \emph{Esta cobertura} tiene como picante algo muy picante.''

%Otro tipo de referencias, se da en el momento de explicar una clase. Cuando se crea una sección hablando acerca de una clase \emph{A}, esa sección no vuelve a crearse si existe una clase \emph{B} cuya descripción requiere que se explique la clase \emph{A}, sino que en la clase \emph{B} se hace referencia a la sección donde ya fue explicada la clase \emph{A}. Esto permite reutilizar las secciones y evita agregar información redundante en el texto.

\subsection{Agregación de sentencias}
La agregación de sentencias ocurre en dos lugares. Uno es durante el proceso de producción de una oración en el constructor \emph{N-Ary}, es decir dentro de una oración; el otro es durante el proceso de creación de párrafos, es decir, entre oraciones.

%de marco metodologico para la gnl
Basaremos la agregación en la conjunción por componentes compartidos~\cite{bernardos2003marco}. El objetivo es que los elementos en común aparezcan una sola vez, mediante elipsis del componente repetido, por ejemplo: sean las oraciones: ``la casa tiene color rojo'', ``la casa tiene color azul'' y ``la casa tiene color verde'', al comparar linealmente las oraciones, se elimina de las consecutivas oraciones la parte inicial que tengan en común. El resultado de esta agregación es: ``la casa tiene color rojo, azul y verde''.
    
\begin{itemize}
    \item {\bf Agregación en constructor NAry}: en las situaciones en las que este constructor se utiliza para enumerar sentencias, como el orden de los elementos enumerados no altera el significado de la oración, se optó por agrupar los elementos según su función (T, SV, ON). Para el grupo de sentencias T, se omitió información inicial de cada sentencia. Por ejemplo: las oraciones ``pizza con carne'', ``pizza picante'' y ``pizza no vegetariana'', se convierten en la oración ``pizza con carne, picante y no vegetariana''.
    
    Para los grupos SV y ON, se ordenaron por verbo, y cada verbo aparece una sola vez, acompañado por la enumeración de los complementos de cada elemento. Por ejemplo: las oraciones ``no tiene como cobertura alguna cobertura de carne'' y ``no tiene como cobertura alguna cobertura de pescado'' se convierten en ``no tiene como cobertura alguna cobertura de carne ni pescado''.
    
    \item {\bf Agregación entre sentencias}: esta agregación se aplica cuando se recorren las sentencias para armar un párrafo. Además de agrupar las sentencias por función (T, SV, ON), también se separan por constructor, obteniendo así 8 grupos: \emph{ObjectSomeValuesFrom, ObjectAllValuesFrom, intersection, union, ObjectHasValue, T, ON y SV}. De esta manera, para cada grupo, se realiza la agregación, de manera independiente entre ellos.
\end{itemize}

\subsection{Ejemplos verbalización de Axiomas}
Con el objetivo de mostrar el proceso de verbalización de forma intuitiva, se representó la composición de oraciones parciales de forma gráfica usando árboles. Se tomaron dos ejemplos de Axiomas de dos ontologías diferentes, uno en lenguaje español y el otro en inglés.

Para no sobrecargar la imagen, se omitió información de menor importancia, como el resultado de verbalizar una Entidad en los nodos hojas, ya que es el mismo nombre de la Entidad pero sin la nomenclatura CamelCase; y se omitió la información gramatical de cada palabra.

\subsubsection{Ejemplo verbalización lenguaje español}
El ejemplo se tomó de la ontología Pizza. El Axioma se corresponde al constructor \emph{owl:equivalentClass} sobre la clase \emph{PizzaPicanteEquivalente}. El contenido del axioma es el siguiente: 
\begin{verbatim}
EquivalentClasses(
    :PizzaPicanteEquivalente
    ObjectIntersectionOf(
        :Pizza
        ObjectSomeValuesFrom(
            :tieneCobertura
            ObjectIntersectionOf(
                :CoberturaDePizza
                ObjectSomeValuesFrom(
                    :tienePicante
                    :MuyPicante)
            )
        )
    )
)    
\end{verbatim}
La verbalización se muestra en la figura~\ref{fig:ejemplo_verb_espaniol}.

Los pasos del proceso se encuentran enumerados desde el 1 al 6. La verbalización comienza en la parte más profunda del árbol, y a medida que se generan las oraciones parciales se retornan al nivel superior del árbol para continuar la composición.

Las composiciones se resuelven a través de las Gramáticas de la sección~\ref{sec:verbalizacion_constructores}. Por ejemplo, en el paso Nº1 utiliza patron$_1$ de la Gramática~\ref{gram:quant_obj_rest}, donde:
\begin{itemize}
    \item $<$verboPropiedad$>$ es ``tiene''.
    \item $<$complementoPropiedad$>$ es ``picante''.
    \item $<$enlace$>$ es ``algún''.
    \item $<$clases$>$ es ``muy picante''.
\end{itemize}

En el paso Nº2, como es una intersección de dos operandos, utiliza la subordinación de SV respecto a T.

De esta manera, en cada nivel busca utilizar un patrón que se ajuste a los componentes recibidos.

\begin{figure}
    \centering
    \includegraphics[width=\textwidth]{img/generacion_documento/verbalizacion_equivalentClass_spanish.pdf}
    \caption{Ejemplo gráfico de verbalización en lenguaje español.}
    \label{fig:ejemplo_verb_espaniol}
\end{figure}

\subsubsection{Axioma lenguaje inglés}
El ejemplo se tomó de la ontología Wine. El Axioma se corresponde al constructor \emph{owl:equivalentClass} sobre la clase \emph{CheninBlanc}. El contenido del axioma es el siguiente: 

\begin{verbatim}
EquivalentClasses(
    :CheninBlanc
    ObjectIntersectionOf(
        :Wine
        ObjectHasValue(
            :madeFromGrape
            :CheninBlancGrape)
        ObjectMaxCardinality(
        1
        :madeFromGrape)
    )
)
\end{verbatim}
La verbalización se muestra en la figura~\ref{fig:ejemplo_verb_ingles}.

Los pasos del proceso se encuentran enumerados desde el 1 al 4. Las gramáticas utilizadas son similares a las explicadas para lenguaje español, con la diferencia que los enlaces han sido traducidos a inglés.

Resulta conveniente explicar el paso Nº3 particularmente. El resultado de este paso es consecuencia de la intersección entre ``Wine'' y el resultado del paso 1 y 2. 
Como es una intersección de tres componentes, se procede a agrupar los componentes según su tipo, para luego unirlos a través de la conjunción ``y''. Dado que el operador \emph{hasValue} y \emph{minCardinality} retornan componentes de tipo SV, son agrupados juntos en una oración. Luego de agruparlos, se procede a realizar el proceso de agregación entre ambas, a través del cual se elimina el verbo ``made'' de la oración parcial construida en el paso 1. 

Por otro lado, la Clase Wine retorna un tipo T, por lo que queda aislado en una oración independiente. Estas agrupaciones resultan en el texto del paso Nº3.

La verbalización de este axioma da como resultado dos oraciones. Como ambas oraciones hacen referencia a la misma Clase, en el paso Nº4 ocurre un reemplazo del nombre de la Clase por una Expresión de Referencia en la segunda oración. 

\begin{figure}
    \centering
    \includegraphics[width=\textwidth]{img/generacion_documento/verbalizacion_equivalentClass_english.pdf}
    \caption{Ejemplo gráfico de verbalización en lenguaje inglés.}
    \label{fig:ejemplo_verb_ingles}
\end{figure}

\subsection{Diagrama de clases}
En la figura \ref{fig:diagrama_clases_microplanificador} se puede ver el diagrama de clases asociado al módulo de Micro Planificación. 

La clase \emph{Statement} es la encargada de realizar la verbalización de los constructores de Axiomas y las enumeraciones. Un objeto \emph{Statement} puede generar más de una oración, ya que cada \emph{Statement} trata un solo tópico y un solo tipo de constructor, pero pueden haber varios axiomas que utilicen el mismo tipo de constructor, por lo que todos esos axiomas estarán en la misma \emph{Statement}.

La clase \emph{StatementComponent} representa la oraciones parciales. Se encarga de componer oraciones parciales y generar una nueva oración parcial. En el caso de un constructor de Entidad, se encarga de generar una oración parcial compuesta por el nombre de la Entidad.

La clase \emph{Word} representa las palabras usadas durante la verbalización. Contiene la función gramatical de la palabra, y los algoritmos necesarios para convertir una palabra en plural, un número en palabra (1 a uno), cambiar el género de una palabra, retornar el artículo correspondiente a un sustantivo, entre otros.

Las subclases de \emph{OWLRestriction} tienen la capacidad de verbalizar los constructores de Expresiones de Clases y de generar las \emph{StatementComponent}.

\begin{figure}
    \centering
    \includegraphics{img/generacion_documento/diagrama_clases_microplanificador.pdf}
    \caption{Diagrama de clases del Micro Planificador.}
    \label{fig:diagrama_clases_microplanificador}
\end{figure}


\section{Implementación Micro Planning}
Para facilitar la implementación de los algoritmos de Agregación y Expresiones de Referencia, decidimos implementar el contenido de las oraciones con listas de palabras en lugar de cadenas de texto, facilitando su manipulación y recorrido. De esta manera, se creó la interfaz \emph{Word} con sus respectivas subclases \emph{WordEnglish} y \emph{WordSpanish}. 
La clase \emph{WordSpanish} utiliza la librería \emph{javaGramatica}\footnote{\url{https://www.proinf.net/permalink/gramatica_numero_genero_y_acentuacion}}, para implementar algunos métodos del idioma español, como pluralizar, cambio de género, número, etc. Para el idioma inglés estos métodos fueron implementados manualmente.

Cada vez que se crea una Word, debe especificarse su palabra y el tipo de función que cumple la palabra. Para identificar el tipo de una palabra, se utilizó el software \texttt{Stanford POS Tagger}\footnote{https://nlp.stanford.edu/software/tagger.html}, el cual permite etiquetar palabras de inglés, español y otros idiomas. Este software utiliza una nomenclatura específica para etiquetar las palabras (llamadas etiquetas \texttt{EAGLE}\footnote{\url{https://web.archive.org/web/20160325024315/http://nlp.lsi.upc.edu/freeling/doc/tagsets/tagset-es.html}} para el español, y para el inglés el \texttt{Penn Treebank}\footnote{\url{https://www.ling.upenn.edu/courses/Fall_2003/ling001/penn_treebank_pos.html}}), por lo que utilizamos esa nomenclatura y la implementamos a través de Variables de Clase creadas en la clase Word.

Los objetos \emph{Word} son creados por cada \emph{StatementComponent}. Una \emph{StatementComponent} se crea durante el proceso de verbalización de los constructores de Expresiones de Clase y de Entidades. Como los constructores del nivel más bajo son de Entidades, es en ellas donde se decide qué tipo de objeto se instancia según el idioma, es decir \emph{StatementComponentEnglish} o \emph{StatementComponentSpanish}. 

En el nivel de los constructores de Expresiones de Clase se encuentran las clases que heredan de la interfaz \emph{OWLRestriction}. Esta interfaz acopla los atributos que requieren los constructores, como los individuos, las clases, las propiedades y cardinalidades.

Las clases que implementan \emph{OWLRestriction} componen y generan las \emph{StatementComponent}. Como ejemplo, se muestra en la Anexo~\ref{sec:clase_OWLObjectComplementOf} la implementación de la clase \emph{OWLObjectComplementOf} con los métodos para implementar la verbalización del constructor en español.

Se puede ver que el método \emph{generateStatementSpanish()} recibe como párametro un objeto de clase \emph{TextCotext}. Esta clase tiene el tópico al que hace referencia la oración que se está componiendo y el lenguaje de la oración, para ser utilizados, por ejemplo, en la generación de Expresiones de Referencia.

Por último, la interfaz \emph{Statement} tiene los atributos necesarios para que sus implementaciones \emph{StatementEnglish} y \emph{StatementSpanish} puedan verbalizar los constructores de Axiomas y de unir todas las oraciones. 

Las \emph{Statement} utilizan la Agregación y las Expresiones de Referencia cuando generan varias oraciones. Estas tareas se implementaron en las clases \emph{StatementComponent} (teniendo en cuenta sus implementaciones) y \emph{OWLClass} respectivamente. En el Anexo~\ref{sec:met_exp_ref} se muestra el método usado para generar las Expresiones de referencia en la clase \emph{OWLClass}.


\section{Diseño del realizador}
%La arquitectura de este trabajo relegó la tarea de definir la estructura de las oraciones y llevar acabo su realización lingüística a la etapa de Micro Planning.
En esta etapa queda la tarea de definir la estructura final del documento. 

Alcanzaremos el Documento Final, a través del refinamiento del Documento Inicial creado en el Macro Planning.

El Documento Inicial tiene una estructura en la que todos los tópicos tienen su propia sección, por lo que no existe más información que haga posible expandir la estructura. Por este motivo, nos enfocaremos en la idea de reducir la estructura del documento convirtiendo secciones en párrafos, para que se haga más compacto, tratando de mejorar la estética y el proceso de lectura, sin perjudicar la coherencia y la segmentación de la información.

Para llevar a cabo el refinamiento, se buscó medir la información a verbalizar, para saber cuándo es posible convertir una sección. Para esto se propuso como métrica principal a la cantidad de oraciones presentes en la descripción de un tópico.

\subsection{Criterio de reducción de Secciones}
Una sección agrupa uno o más párrafos o sub-secciones, por lo que en una sección puede existir más de un tópico, con la condición de que esos tópicos tengan una relación, ya sea de hermanos o de hijos. Estas relaciones se obtienen del Documento Inicial. 

Para decidir si un tópico es planificado como una Sección, se utilizan los siguientes criterios:
\begin{itemize}
    \item Debe poseer cinco o más oraciones.
    \item Debe poseer alguna Subsección con cinco o más oraciones.
    \item En cualquier otro caso será planificado como un Párrafo.
\end{itemize}

Si la cantidad de oraciones requeridas para ser considerado una sección, tiende a cero, la planificación del documento será similar al Documento Inicial; mientras que si la cantidad de oraciones tiende a infinito, la planificación perderá niveles de jerarquía, y únicamente contemplará como secciones a los tópicos del primer nivel. Ninguno de los dos extremos es conveniente: o se pierde coherencia, o se obtiene un texto con demasiados niveles de secciones, que resulta antinatural. Sin embargo, un documento en el que predominan las sub-secciones (aunque tengan poca información), ayuda a los humanos a reconocer y establecer relaciones entre los tópicos, siendo el objetivo principal no desaprovechar la semántica del dominio modelado, y teniendo como intención principal que el lector comprenda el dominio modelado en la ontología.

\section{Implementación del Realizador}
El Realizador está embebido dentro de la clase Section, en forma de métodos que implementan los criterios vistos para la definición del documento. En el Anexo~\ref{sec:clase_section} se presenta un pseudocódigo de la clase Section.

\section{Casos de estudio}
En el capítulo de Organización de Información analizamos cómo quedan organizados los tópicos de las ontologías Pizza y Wine, por lo que continuaremos con estas ontologías de ejemplo, analizando la realización del texto generado. Para identificar el formato del texto, a las secciones se le incluye el número de sección a la izquierda del título de la sección, y cada párrafo comienza con sangría.

Ya que existen muchos parámetros a tener en cuenta en el análisis de un texto, tendremos en cuenta criterios que estén directamente relacionados con el Objetivo propuesto: \begin{itemize}
    \item En cuanto a la estructura jerárquica del texto, analizaremos la coherencia global y si los temas son representativos de la ontología.
    \item Respecto a la sintaxis, veremos si las oraciones son claras, cumplen con la sintaxis del idioma, son coherentes a nivel local, y si transmiten el significado del axioma verbalizado. 
    \item En cuanto a la realización del texto, analizaremos la estética (aunque de manera subjetiva), viendo si la segmentación en secciones y párrafos es agradable.
\end{itemize}


\subsection{Generación documento ontología Pizza}
Ya que el texto completo resulta muy extenso, solo pondremos los fragmentos más significativos del Documento Final. En la figura~\ref{fig:doc_final_pizza} se presentan los fragmentos a analizar. Para visualizar la estructura total del documento, se imprimió el índice de las secciones, el cual se muestra en la Figura~\ref{fig:indice_secciones_pizza}. 

Los tópicos reconocidos representan bien el tema del texto. Se puede apreciar que la estructura del documento comparte algunas similitudes con el Documento Inicial, manteniendo la coherencia global, pero disminuyó la cantidad de secciones, ya que muchas de ellas se convirtieron en párrafos. Aún así, prevalecen títulos de secciones que agregan información semántica aunque no aporten información en forma de texto, tales como Ingredientes, Coberturas y Bases.

La sintaxis es aceptable y las oraciones pueden comprenderse, aunque pueden mejorarse para disminuir la redundancia. Los axiomas no son muy complejos, aunque verbalizó bien axiomas de hasta cuatro niveles de jerarquía, como el de la Figura~\ref{fig:ejemplo_verb_espaniol}. Las oraciones resultan coherentes, el salto entre oraciones es fluido teniendo en cuenta que no intercalan tópicos, y el contenido del texto anticipa bien la aparición de párrafos y subsecciones.

La estructura de las oraciones es repetitiva, no era objetivo de este trabajo realizar un texto variado, por lo que ante la misma información para verbalizar se obtiene la misma estructura.

En cuanto a la cuestión estética, el texto resulta agradable, no posee muchas secciones anidadas por lo que visualmente no sobrecarga la lectura.

\begin{figure}
\centering
\fbox{
\begin{minipage}{14cm}
\setlength{\parindent}{1em}
\large{
\noindent 1 Pizza
}

\small{
Una pizza es una comida. Una pizza tiene base de pizza y tiene cobertura de pizza. Esta tiene como base a alguna base de pizza. Existen las siguientes clases de pizzas: pizza con carne, picante, no vegetariana, con nombre, \dots, con queso y vegetariana. Una pizza es disjunta de base de pizza, cobertura de pizza y helado.

Pizza con carne: una pizza con carne es una pizza que tiene como cobertura a alguna cobertura de carne. 

Pizza vegetariana equivalente1: una pizza vegetariana equivalente1 es una pizza que tiene como cobertura exclusivamente a cobertura vegetariana.

Pizza picante equivalente: una pizza picante equivalente es una pizza que tiene como cobertura a alguna cobertura de pizza que tiene como picante algo muy picante. 

Pizza real italiana: una pizza real italiana es una pizza que tiene pais de origen italy. Esta tiene como base exclusivamente a base delgada y crujiente.

\par \dots Otras pizzas realizadas como párrafos.
}

\large{
\noindent 1.1 Pizza con nombre
}

\small{
Una pizza con nombre es una pizza. Existen las siguientes clases de pizzas con nombre: margherita, frutti di mare, \dots, veneziana y napoletana.

Margherita: una margherita es una pizza con nombre. Esta tiene como cobertura a alguna cobertura de mozzarella y tomate. También, margherita tiene exclusivamente las siguientes coberturas: cobertura de mozzarella o de tomate. una margherita es disjunta de american, american hot, \dots, soho y veneziana.

\dots Otras pizzas con nombre realizadas como párrafos.
}

\large{
\noindent 1.2 Ingredientes
	
\noindent 1.2.1 Coberturas
	
\noindent 1.2.1.1 Cobertura de pizza
}

\small{
Una cobertura de pizza es una comida. Una cobertura de pizza es cobertura de una pizza. Existen las siguientes clases de coberturas de pizza: \dots
}

\large{
\noindent 1.2.2 Bases
}

\large{
\noindent 1.2.2.1 Base de pizza
}

\small{
Una base de pizza es una comida. Una base de pizza es base de una pizza. Existen las siguientes clases de bases de pizza: base gruesa y delgada y crujiente.
}
\end{minipage}
}
\caption{Fragmentos del Documento Final generado a partir de la ontología Pizza.}
\label{fig:doc_final_pizza}
\end{figure}

\begin{figure}
\begin{multicols}{2}
\small{
\begin{figure}[H]
\dirtree{%
.1 Pizza.
.2 Pizza con nombre.
.2 Ingredientes.
.3 Coberturas.
.4 Cobertura de pizza.
.5 Cobertura de verduras.
.6 Cobertura de pimiento.
.6 Cobertura de tomate.
.5 Cobertura de hierbas.
.5 Cobertura de carne.
.5 Cobertura de queso.
.5 Cobertura de pescado.
.3 Bases.
.4 Base de pizza.
.5 Base gruesa.
.5 Base delgada y crujiente.
}
\end{figure}

\begin{figure}[H]
\dirtree{%
.1 Otras secciones.
.2 Domain concept.
.3 Comida.
.2 Value partition.
.3 Picante.
}
\end{figure}
}
\end{multicols}
\caption{Índice de secciones del Documento Final generado para la ontología Pizza.}
\label{fig:indice_secciones_pizza}
\end{figure}


\subsection{Generación documento ontología Wine}
Al igual que el caso de estudio anterior, solo incluiremos algunos fragmentos del documento generado, los cuales se pueden ver en la Figura~\ref{fig:doc_final_wine}.

En la Figura~\ref{fig:indice_secciones_wine} se visualiza el índice del documento. El resultado es similar al obtenido en la verbalización de la ontología Pizza, comparando la Figura~\ref{fig:indice_secciones_wine} con la Figura~\ref{fig:caso_estudio_wine}, se observa que la estructura resulta mucho más compacta.

Analizando la coherencia global y el sentido del texto, la estructuración resultó razonable, enfocada en los vinos y las propiedades de los vinos.

La implementación de la sintaxis es sustancialmente equivalente a la del español, por lo que puede presentar las mismas características y capacidades, como redundancia, repetitividad y claridad. 

El resultado de la realización final también resulta satisfactorio. Si bien presenta hasta cinco niveles de anidamiento (a partir de \emph{Consumable thing}), la información que se transmite es necesaria para mantener la coherencia global y establecer las relaciones pertinentes entre los tópicos.

\begin{figure}
\centering
\fbox{
\begin{minipage}{14cm}
\setlength{\parindent}{1em}
\large{
\noindent 1 Wine
}

\small{
A wine is a potable liquid. A wine has wine sugar as sugar, has wine flavor as flavor, has wine body as body, has wine color as color, has wine descriptor and made from grape a wine grape. This located in some region. Also, wine has only winery maker. Also, wine has exactly one maker, made at least one wine grape, wine body, wine color, wine flavor and wine sugar. There are the following kinds of wines: italian wine, \dots, late harvest and alsatian wine. 

Full bodied wine: a full bodied wine is a wine that has full body.

Burgundy: a burgundy is a wine that located in bourgogne region. This has dry sugar. There are the following kinds of burgundies: white burgundy and red burgundy. 

\dots Otros vinos realizados como párrafos.
}

\large{
\noindent 1.1 Italian wine
}

\small{
A italian wine is a wine that located in italian region. chianti is the only kind of wine. 

Chianti: a chianti is a italian wine. This has only light or medium body. Also, chianti has red color, moderate flavor and dry sugar. located in chianti region. made from grape sangiovese grape. chianti classico is a type of chianti. 

Chianti classico: a chianti classico has mc guinnesso maker and medium body.
}

\large{
\noindent 1.2 ... 1.7
}

\large{
\noindent 1.8 Red wine
}

\small{
A red wine is a wine that has red color. There are the following kinds of wines: dry red wine, red burgundy, port and red bordeaux. 

Port: a port is a red wine. This has full body, strong flavor and sweet sugar. located in portugal region. taylor port is a type of port. 
Taylor port: a taylor port has taylor maker. 
}


\large{
\noindent 1.8.1 Red burgundy
}

\small{
A red burgundy is a burgundy and red wine. This made from grape pinot noir grape. Also, red burgundy made at most one wine grape. cotes d or is the only kind of burgundy. 
Cotes d or: a cotes d is a red burgundy that located in cotes d or region. This has moderate flavor. clos de vougeot cotes d or is a type of cotes d or. 
Clos de vougeot cotes d or:	a has clos de vougeot maker. 
}

\small{\dots Otros vinos realizados como secciones}

\large{
\noindent 1.14 Wines descriptor
}

\small{
Wine sugar:	a wine sugar is a wine taste. a wine sugar is a sweet, off dry or dry. there are the wines sugar: dry, off dry and sweet. 

Wine color: a wine color is a wine descriptor. a wine color is a rose, red or white. there are the wines color: white, red and rose. 

Wine flavor: a wine flavor is a wine taste. a wine flavor is a moderate, delicate or strong. there are the wines flavor: moderate, strong and delicate. 

Wine body: a wine body is a wine taste. a wine body is a light, medium or full. there are the wines body: medium, full and light. 
}

\large{
\noindent Otras secciones menos importantes
}

\large{
\noindent 2 Region
}

\small{
A region adjacent region. There are the regions: arroyo grande region, california region, ....

Arroyo grande region: an arroyo grande region located in california region. 

....
}

\large{
\noindent 3 Vintage year
}

\small{
A vintage year has year value. year1998 is a type of year. 
Year1998: a year1998 has year value. a year1998 has year value. 
}

\small{
\dots Otras secciones.
}

\end{minipage}
}
\caption{Fragmentos del Documento Final generado a partir de la ontología Wine.}
\label{fig:doc_final_wine}
\end{figure}


\begin{figure}
\begin{multicols}{2}
{\small
\begin{figure}[H]
\dirtree{%
.1 Wine.
.2 Italian wine.
.2 Loire.
.2 Table wine.
.2 Pinot noir.
.2 Dessert wine.
.3 Sweet riesling.
.2 Bordeaux.
.3 Medoc.
.2 Riesling.
.2 Red wine.
.3 Red burgundy.
.2 Cabernet sauvignon.
.2 Semillon or sauvignon blanc.
.3 Sauvignon blanc.
.2 White wine.
.3 White burgundy.
.2 Zinfandel.
.2 Chardonnay.
.2 Wines descriptor.
.3 Wine sugar.
.3 Wine color.
.3 Wine flavor.
.3 Wine body.
}
\end{figure}}  

\begin{figure}[H]
\dirtree{%
.1 Otras secciones.
.1 Region.
.1 Vintage year.
.1 Wine descriptor.
.1 Non consumable thing.
.1 Vintage.
.1 Fruit.
.1 Winery.
.1 Consumable thing.
.2 Meal course.
.3 Foods.
.4 Edible thing.
.5 Seafood.
.5 Pasta.
}
\end{figure}

\end{multicols}
\caption{Índice de secciones del Documento Final generado para la ontología Wine.}
\label{fig:indice_secciones_wine}
\end{figure}


Ya que en lenguaje inglés hay ejemplos de verbalizaciones de otras investigaciones, a continuación se presenta una comparación de la verbalización de la clase \emph{Beaujolais} de la ontología \emph{Wine}, entre el trabajo~\cite{hewlett2005effective} y el nuestro:

En~\cite{hewlett2005effective} proponen listar sus características de la siguiente manera: 
``A Beaujolais is a Wine that:
\begin{itemize}
    \item is made from at most 1 grape, which is Gamay Grape
    \item has Delicate flavor
    \item has Dry sugar 
    \item has Red color
    \item has Light body ''
\end{itemize}

En nuestro caso, la oración resultante es la siguiente:
``Beaujolais has light body, red color, delicate flavor and dry sugar. made from grape gamay grape.''

Existe clara diferencia, no solo en el criterio de presentación lista frente a prosa, sino también en el hecho de elegir usar o no elipsis (en nuestro caso omitimos el verbo ``has'' para evitar redundancia), y de separar en diferentes oraciones las características asociadas al ``has'' y las asociadas al ``made from''. Decidir qué representación usar depende de muchos factores, por ejemplo, la versión en prosa no resulta muy compleja, pues la omisión del ``has'' solo ocurre en tres características, por lo que no recae tanta carga sobre el lector. Quizá en descripciones más largas sea conveniente un formato en lista, pero también hay que tener en cuenta cómo impacta visualmente representar todas las clases en forma de listas en una gran ontología. De acuerdo con~\cite{hewlett2005effective}, pueden existir métodos de generación más complejos que mejoren la calidad del texto de salida.

Otra herramienta que se pudo probar con la ontología \emph{Wine} fue Attempto Tools, la cual retornó la siguiente salida para la clase Beaujolais:

``
Every Beaujolais hasBody Light.

Every Beaujolais hasColor Red.

Every Beaujolais hasFlavor Delicate.

Every Beaujolais hasSugar Dry.

Every Beaujolais madeFromGrape GamayGrape.

Every Beaujolais madeFromGrape at most 1 thing.
''

En comparación con las salidas en formato lista y prosa, podemos apreciar que esta salida es la que menos trabaja sobre las oraciones, repitiendo información y dejando la nomenclatura camelCase sin procesar.

\section{Conclusiones}
Hemos presentado un sistema que genera un texto en lenguaje natural a partir de los axiomas de una ontología. La estructura del documento de texto depende de la organización realizada en el macroplaning, y de los criterios tomados en el Realizador, conservando la coherencia global recibida del Organizador de Información. 

Tomando como referencia lo postulado por la Teoría de Centrado, se pudo proveer a los párrafos de coherencia local. Para cada clase se reunió la mayor cantidad de información disponible, agrupándola a través de los constructores de Axiomas.

La separación en constructores de Expresiones de Clases y constructores de Axiomas, dio lugar a las oraciones parciales y a un sistema de tipos. Si bien este sistema no se corresponde fielmente a las categorías de constituyentes sintácticos\footnote{Un constituyente sintáctico es una palabra, o secuencia de palabras, que funciona en conjunto como una unidad dentro de la estructura jerárquica de una oración} de alguna gramática conocida (como la gramática tradicional o la gramática generativista), permitió minimizar la cantidad de combinaciones que había que diseñar en las gramáticas.

Las sentencias fueron debidamente procesadas para aplicar algunas técnicas de Procesamiento de Lenguaje Natural, como expresiones de referencia, agregación de sentencias, inflexión de número y determinación de pronombres basados en el género (en el caso del lenguaje español). 

El resultado en los casos de prueba fue satisfactorio, la sintaxis generó oraciones interpretables.

Como mostramos en las comparaciones con resultados de otros trabajos, es posible tomar diferentes criterios para mostrar la misma información, y desarrollar técnicas o gramáticas más sofisticadas para mejorar la fluidez del texto.

\chapter{Conclusiones}
En el desarrollo de esta tesis se diseñó e implementó un Sistema de Generación de Lenguaje Natural, orientado a generar un texto organizado a partir de una ontología (OWL) de la Web Semántica. Para desarrollar tal sistema se propusieron dos problemas fundamentales: el de organización de la información y el de generar el documento de texto. 

%Para abordar el problema de la organización de la información, se tomó en cuenta que los textos en general  poseen una estructura jerárquica en el que sus tópicos presentes están semánticamente relacionados, al igual que las ontologías. Para maximizar la relación semántica entre los tópicos, se usó como base la coherencia global, aplicando los fundamentos de la Teoría de Veins. Se propuso una solución basada en el análisis del grafo subyacente a la ontología, midiendo las relaciones entre sus nodos. De esta manera se obtuvieron las Entidades de la ontología que resultaban más sobresalientes, con el fin de usarlas como hilo conductoras de los temas del texto.
Para abordar el problema de la organización de la información, se tomó en cuenta que cualquier texto  posee una estructura jerárquica en el que sus tópicos  están semánticamente relacionados. Esta característica no queda exenta en las ontologías en donde también podemos evidenciar este comportamiento. 

Para maximizar la relación semántica entre los tópicos, se usó como base la coherencia global, aplicando los fundamentos de la Teoría de Veins. Se propuso una solución basada en el análisis del grafo subyacente a la ontología, midiendo las relaciones entre sus nodos. De esta manera se obtuvieron las Entidades de la ontología que resultaban más sobresalientes, con el fin de usarlas como hilo conductor de los temas del texto generado.
Esta solución se vio motivada por la característica que poseen las ontologías de la Web Semántica, de ser grafos libres de escala. 

% me hace ruido esta oracion por el "satisfactorio" dado que hay que evitar comentarios o auto evaluaciones en esta parte. Aquí se centra la atención  los conocimientos que la tesis propone y en este caso en los casos de estudios,.
% Por ahi habría que darla vuelta y decir que la ve un mayor imparcto en el texto resultante cuando la cant de objectPropties y del contenido de sus dominios .. Sin embargo .... (Sin la valoridad)

Los resultados obtenidos de los casos de prueba retornaron estructuras organizadas, que se correspondían con los resultados esperados. Sin embargo, la solución propuesta depende de la cantidad de \emph{ObjectProperties} y del contenido de sus dominios, por lo que queda pendiente poder analizar la solución con mayor diversidad de ontologías, variando la densidad de sus relaciones.

Además de maximizar la relación entre los tópicos, se buscó maximizar la relación entre los elementos de los párrafos. Para esto se usó como base la coherencia local, a través de los postulados de la Teoría de Centrado.

Para el sistema de generación de texto, se dividió el problema en tres sub-problemas, basados en los procesos conocidos del Procesamiento del Lenguaje Natural. Estos subproblemas fueron diseñar e implementar la Macroplanificación, Microplanificación y la Realización. 

Para comenzar, se desarrolló el módulo Macroplanificador, encargado de brindar la estructura jerárquica del texto. Dado que se integró el Macroplanificador con el Organizador de Información, su implementación fue trivial.
Luego se procedió a verbalizar los axiomas de la ontología en la Microplanificación, desarrollando la gramática del lenguaje humano e implementando los algoritmos de generación de texto adecuados, para reducir la redundancia y mejorar la fluidez del texto.
Por  último  se  resolvió  el  problema  de  la  Realización  del  texto,  modificando  la  estructura inicial obtenida del Macroplanificador, para reducir los niveles de la jerarquía del texto donde era adecuado.

Como se mostró en los casos de prueba, la solución basada en el desarrollo de estos módulos resultó en textos legibles y organizados. Si bien validar textos puede resultar subjetivo, podemos apreciar que existe suficiente relación semántica en el orden de los temas que se presentan, tanto entre las Secciones como entre las Oraciones, por lo que cumplen con cierta Coherencia Global y Local.

En cuanto a la sintaxis, se demostró que pueden existir mejores formas de presentar la información, que depende de cada caso particular.

Aún resulta conveniente explorar otras técnicas para desarrollar las gramáticas y la composición de las oraciones. Desarrollar las gramáticas manualmente es un trabajo costoso, y debido a la naturaleza recursiva de las construcciones de los axiomas y de la semántica de sus constructores, es difícil mantener el control de la oración compuesta, tanto a nivel sintáctico como semántico. Queda pendiente probar el sistema con otras ontologías que presenten mayor complejidad en sus axiomas, con el fin de desarrollar nuevas gramáticas que abarquen más casos de prueba.

Más allá de probar el sistema con nuevos casos de prueba, para reconocer las ventajas y desventajas de la solución propuesta, aún hay ciertos factores a determinar, para ayudar a mejorar la estética del texto, que quedan fuera del alcance de esta tesis, por ejemplo:
\begin{itemize}
    \item Decidir cuándo es conveniente utilizar texto en formato lista en lugar de prosa, o en algún otro formato como tablas comparativas.
    \item Decidir cuándo es conveniente utilizar expresiones de referencia o elipsis.
    \item Establecer criterios más sofisticados para modificar la estructura del texto en la etapa de Realización. 
\end{itemize}


\bibliographystyle{abbrv}
\bibliography{bibliografia}

\chapter*{Apéndice}
\appendix
%\markboth{Appendices}{}
\addcontentsline{toc}{chapter}{Apéndice}
\renewcommand{\thesection}{A.\arabic{section}}

\section{Implementación Organizador de Información}
\label{apx:impl_org_inf}

\subsection{Clase \texttt{OntologyManager}}
\label{sec:clase_ontologymanager}

\begin{minted}[autogobble,linenos]{java}
public class OntologyManager {

    /*nameOf guarda el valor 0 o 1 para reconocer de donde tiene que
    extraer el nombre de las entidades*/
    private static int nameOf; // 0 -> iri, 1 -> label.
    /*lang setea el valor del lenguaje para buscar los nombres
    de las entidades en los labels correctos*/
    private static String lang; // "en" -> ingles, "es" -> español.

  public static Ontology createOntology
            (OWLOntology ontology, int nameOfClass, String language) {
        nameOf = nameOfClass;
        lang = language;
        Ontology onto = new Ontology();
        
        //Obtiene el título de la ontología desde la etiqueta title.
        String title = getTitleOntology();
        onto.setTitle(title);
        
        //crear clases
        OWLClass rootClass = createClasses();

        //crear los dataProperty
        OWLDataProperty rootDataProp = createDataProperty();
        
        //crear properties
        OWLObjectProperty rootProp = createObjectProperty();
        //setea las objectProperty equivalentes, inversas y disjuntas.
        setListsToProperties();
        //setea el dominio y rango de las objectProperty
        setObjectPropertyDomainRange();
        //setea el dominio y rango de las dataProperty
        setDataPropertyDomain();
        
        //crear individuals
        createIndividuals();
        
        /*crear las clases anónimas para cada clase en onto.
        Setea listas para guardar los axiomas. Para cada clase en onto
        setea sus clases equivalentes, disjuntas y superClases.*/
        /*createAnonymousClasses debe ser llamado despues de haber
        creado las clases nombradas, las properties y los individuos.*/
        createAnonymousClasses();
        
        //agrega los componentes creados al objeto onto.
        addComponentsToOnto();
        
        //precomputeValues asocia a cada clase en onto una lista 
        //de propiedades en las que participa como dominio.
        onto.precomputeValues();
        return onto;
    }
}
\end{minted}

\subsection{Clase \texttt{Ontology}}
\label{sec:clase_ontology}
\begin{minted}[autogobble,linenos]{java}
public class Ontology {

    private String title;
    //guardan las Entidades de nivel mas alto en la jerarquía.
    private OWLClass rootClass;
    private OWLObjectProperty rootObjectProperty;
    private OWLDataProperty rootDataProp;

    //los siguientes hashmaps asocian IRIs a sus respectivas Entidades.
    private HashMap<IRI, OWLClass> classes;
    private HashMap<IRI, OWLObjectProperty> properties;
    private HashMap<IRI, OWLObjectProperty> anonProperties;
    private HashMap<IRI, OWLDataProperty> dataProperties;
    private HashMap<IRI, OWLIndividual> individuals;
    
    //Las siguientes listas guardan las Entidades presentes en la 
    //ontologia
    private LinkedList<OWLClass> listClasses;
    private LinkedList<OWLObjectProperty> listProps;
    private LinkedList<OWLObjectProperty> anonListProps;
    private LinkedList<OWLDataProperty> listDataProps;
    private LinkedList<OWLIndividual> listIndividuals;

    /*hashmaps para guardar que clases participan en el dominio
    de que propiedades*/
    HashMap<OWLClass, LinkedList<OWLObjectProperty>> classOnPropDomain;
    /*hashmaps para guardar que clases participan en el rango
    de que propiedades*/
    HashMap<OWLClass, LinkedList<OWLObjectProperty>> classOnPropRange;

    /*Retorna una lista de ObjectProperties que sean superProperties
    en comun de las properties de la lista prop.*/
    public LinkedList<OWLObjectProperty>
    getCommonSuperProperties(LinkedList<OWLObjectProperty> prop) {}

    /*Dada la clase cls, recorre las properties donde participa como dominio
    y busca las superProperties de mas alto nivel que encuentre 
    y las retorna. */
     public LinkedList<OWLObjectProperty>
     getCommonSuperPropertiesFromClass(OWLClass cls) {}

    /*Retorna el contenido necesario para crear un nuevo grupo
    de entidades a partir de la clase cls. Util para
    el modulo ContentGrouping */
    public LinkedList<OWLIntClass> 
    createTopicsGroupsFromClass(OWLClass cls) {
        LinkedList<OWLIntClass> res = new LinkedList<>();
        res.addAll(getCommonSuperPropertiesFromClass(cls));
        res.addAll(cls.getSubClass());
        res.addAll(cls.getIndividuals());
        return res;
    }
}
\end{minted}


\subsection{Clase \texttt{ContentClasification}}
\label{sec:clase_content_clasif}

\begin{minted}[autogobble,linenos]{java}
public class ContentClasification {

    /* Es el método principal que retorna las entidades
    más relevantes para agregar en el primer nivel de
    la jerarquía del Árbol de Entidades. */
    public static HashMap<OWLIntClass, LinkedList<OWLObjectProperty>>
    getMainContent(Ontology ontology) {
        
        cls2Cant = getClassesToCantProperties(ontology);
        mainClasses = getMainClasses(cls2Cant);

        /*insertar en un nuevo hashmap cls2SuperProp las clases de mainClasses
        asociadas a las superProperties en común que haya entre las properties 
        donde participa como dominio*/
        for (OWLClass mainClass : mainClasses) {
            cls2SuperProp.put(mainClass,
            ontology.getCommonSuperPropertiesFromClass(mainClass));
        }
        
        /*Si existe una Clase c y alguna SuperClase de c, 
        entonces c se elimina por ser mas específica*/
        cls2SuperProp = removeSubClass(cls2SuperProp);

        /*si el hashmap es vacio, entonces se agregan todas las
        clases sin objproperties*/
        if (cls2SuperProp.isEmpty()) {
            for (OWLClass cls : ontology.getClasses()) {
                cls2SuperProp.put(cls, new LinkedList<>());
            }
        }
        return cls2SuperProp;
    }
    
    /*Devuelve una Lista de HashMaps desde Clases a Integers. Cada
     Hashmap tiene una sola entrada, que va desde una Clase a la
     cantidad de properties en las que participa como dominio.*/
    private static LinkedList<HashMap<OWLClass, Integer>>
    getClassesToCantProperties(Ontology ontology) {}
    
    /*Devuelve una lista de Clases. Las Clases retornadas son las
     que participan en mas cantidad de properties, y que superan 
     el promedio calculado como la suma de todas las properties
     dividido entre la cantidad de clases.*/
    private static LinkedList<OWLClass>
    getMainClasses(LinkedList<HashMap<OWLClass, Integer>> cls) {}
}

\end{minted}


\subsection{Clase \texttt{ContentGrouping}}
\label{sec:clase_content_grouping}

\begin{minted}[autogobble,linenos]{java}
public class ContentGrouping {

    /*A partir de la clase cls y la rama branch a la que pertenece,
    se genera un nuevo grupo de Entidades, y se verifica que cada
    Entidad aún no pertenezca a la rama*/
    public static LinkedList<OWLIntClass>
    generateNewTopics(Ontology ontology, OWLClass cls, 
                LinkedList<OWLIntClass> branch) {
        LinkedList<OWLIntClass> tops = ontology.createTopicsGroupsFromClass(cls);
        for (OWLIntClass r : tops) {
            if (!r.getType().equals("property")) {
                    if (!branch.contains(r)) {
                        res.add(r);
                    }
            }
        }
        for (OWLIntClass r : tops) {
            if (r.getType().equals("property")) {
                checkProperties((OWLObjectProperty) r, res, branch);
            }
        }
        return res;
    }

    /*A partir de una property prop y la clase dom que pertenece
    a su dominio, busca crear un nuevo grupo de Entidades.
    Primero busca subProperties de prop. Si no tiene subProperties, 
    entonces intenta agregar el Rango de prop.*/
    public static LinkedList<OWLIntClass> 
    generateNewTopics(Ontology ontology, OWLObjectProperty prop,
                    OWLClass dom) {
        
        /*obtiene las subProperties de prop*/
        LinkedList<OWLIntClass> newTopics = prop.getSubPropWithDomain(dom);
        
        if (newTopics.isEmpty()) {
            newTopics = prop.getRange();
        }
        return newTopics;
    }
    
    /*A partir de prop busca agregar a ella o una superProperty
    de ella en res. Agrega a prop si tiene una superProperty en branch.
    Si no puede agregar a prop busca agregar un ancestro de prop que 
    tenga una superProperty en branch.*/
    private static boolean checkProperties(OWLObjectProperty prop,
    LinkedList<OWLIntClass> res, LinkedList<OWLIntClass> branch) {}
}
\end{minted}

\subsection{Clase \texttt{GeneratorTreeManager}}
\label{sec:clase_generator_tree}

\begin{minted}[autogobble,linenos]{java}
public class GeneratorTreeManager {

    public static LinkedList<OWLIntClass> getText(Ontology onto) {

        HashMap<OWLIntClass, LinkedList<OWLObjectProperty>> mainClass =
        ContentClasification.getMainContent(onto);
        
        /*agregar las clases de mainClass al primer nivel del árbol. 
        mainEntities representa el primer nivel del árbol 
        (las entidades mas relevantes)*/
        mainEntities.add(mainClass);
        
        /*para cada clase de mainClass crear una nueva rama y agregar 
        sus Nuevos grupos de Entidades*/
        for (Map.Entry<OWLIntClass, LinkedList<OWLObjectProperty>> entrySet : 
                mainClass.entrySet()) {
            /*navigation mantiene las clases del primer nivel y
            las Entidades de la rama que se está creando*/
            navigation.add(mainClass);
            OWLClass mClass = entrySet.getKey();
            LinkedList<OWLIntClass> newEntities =
            ContentGrouping.generateNewTopics(onto, mClass, null);
            addBranchFromEntity(newBranch, mClass, newEntities, navigation,
            onto);
            mainEntities.addNewBranchToEntity(mClass, newBranch);
            navigation.clear();
            
        }
        /*agregar informacion marginada*/
        mainEntities.addNewBranch(addAbsentClass(onto));
        return mainEntities;
    }
    
    /*Recorre las newEntities creando sus respectivos niveles de jerarquía 
    en el Árbol de Entidades.*/
    private static void addBranchFromEntity(newBranch, OWLClass actualClass,
    newEntities, navigation, onto) {
        for (OWLIntClass enti : newEntities) {
            if (!navigation.contains(enti)) {
                navigation.addLast(enti);
                if (enti.getType().equals("class")) {
                    sub2 = ContentGrouping.generateNewTopics(onto, enti, navigation);
                    addBranchFromEntity(newSubBranch, enti, sub2, navigation, onto);
                }else if (enti.getType().equals("property")) {
                    sub2 = ContentGrouping.generateNewTopics(onto, enti, actualClass);
                    addBranchFromEntity(newSubBranch, actualClass, sub2, navigation, onto);
                }
                navigation.removeLast();
                newBranch.addNewBranchToEntity(enti, newSubBranch);
            }
        }
    }
}

\end{minted}


\end{document}
